% Definice vzhledu a nastavení se nachází přímo v následující speciální třídě
\documentclass{template-definitions/thesistbuinzlin}

% Uživatelské definice -- upravte dle potřeby
\setthesislanguage{CZ}
% CZ   -- práce bude v českém jazyce
% EN   -- práce bude v anglickém jazyce

\setfaculty{FAI}
  % FAI  -- pro Fakultu aplikované informatiky
  % FAME -- pro Fakultu managementu a ekonomiky
  % FHS  -- pro Fakultu humanitních studií
  % FLKR -- pro Fakultu logistiky a krizového řízení
  % FMK  -- pro Fakutlu mutimediálních komunikací
  % FT   -- pro Fakultu technologickou
  % UNI  -- pro Univerzitní institut

\setthesistype{DP}
  % BP   -- bakalářská práce
  % DP   -- diplomová práce

\setthesisyear{2026}
  % zadejte rok namísto "xxxx"

\setcitationmode{numbers}
  % authoryear -- odkazování na zdroje bude v módu (autor, rok), tzv. harvardský systém, vhodné v případě fakult FAME, FHS, FLKŘ, FMK, FT
  % numbers    -- odkazování na zdroje bude v číselném módu [číslo], vhodné v případě fakult FAI, FT

% Lze přidat vertikální odsazení nad (prvni parametr) a pod (druhý parametr) obrázky, tabulky i rovnice/soustavy rovnic
\setmarginaroundfigures{0mm}{0mm}
\setmarginaroundtables{0mm}{0mm}
\setmarginaroundequations{0mm}{0mm}

\setauthor{Bc. Jozef Hreha}

\settitlecz{Aplikace pro podporu správy repozitáře vědeckých výstupů za pomocí AI (LLM)}
\settitleen{Thesis Title in English (max. 7 rows)}  % Needed only when writing in English

\setabstractcz{Text abstraktu v jazyce práce.}

\setabstracten{Text abstraktu ve světovém jazyce (angličtině).}

\setkeywordscz{klíčové slovo, klíčové slovo}

\setkeywordsen{keyword, keyword}

\setuppdfa

%\usecolouredlinks  % Volitelné, může být nápomocné při tvorbě a kontrole výsledného dokumentu

% ============================================================================ %

\begin{document}

\maketitle

% nasledujuci riadok tu musi byt, inak sa dobabrú riadky
\insertassignment{graphics/placeholder-scan.pdf}

\insertdeclaration{CZ}
  % CZ   -- doporučované, i pokud práci píše český/slovenský student v angličtině
  % EN   -- variant only for international students

\insertabstractandkeywords

% ============================================================================ %

\clearpage
Zde je místo pro případné poděkování, popř. motto, úryvky knih atp.
Pokud poděkování nechcete vložit, smažte celý tento blok 
(mezi příkazy cleardoublepage, včetně těchto příkazů).



% ============================================================================ %

\insertcontents  % Obsah je generován automaticky

% ============================================================================ %

%\insertlistoffigures  % Seznam je generován automaticky

%\insertlistoftables  % Seznam je generován automaticky

%%\insertlistoflisting  % Seznam je generován automaticky, vhodné jen pokud máte v dokumentu fragmenty kódu

%% ============================================================================ %
% Encoding: UTF-8 (žluťoučký kůň úpěl ďábelšké ódy)
% ============================================================================ %

% It is also possible to automate the generation of this list. Inspirative packages: acro or glossaries.

\inserttitlelistofabbreviations

{\hskip-3mm\setstretch{1.15}
\begin{tabular}{>{\rmfamily\normalsize}p{20mm} >{\rmfamily\normalsize}p{127mm}} % Lze měnit šířku sloupců, součet však musí dávat 147 mm

CPU & Central Processing Unit \\
PTFE & Polytetrafluoroethylene \\
VNA & Vector Network Analyser \\

\end{tabular}
}

% ============================================================================ %
  % Seznam zkratek

%\insertlistofappendices  % Seznam příloh je generován automaticky

% ============================================================================ %

\paragraphindentationstart  % Nastaví odsazování odstavců dle zvoleného jazyka
%\cleardoublepage					  % Další text (úvod) až na liché straně

% Text práce
% ============================================================================ %
% Encoding: UTF-8 (žluťoučký kůň úpěl ďábelšké ódy)
% ============================================================================ %

\sectionwithoutnumbering{Úvod}

První řádek prvního odstavce v kapitole či podkapitole se neodsazuje, ostatní ano. Vertikální odsazení mezi odstavci je typické pro anglickou sazbu; czech babel toto respektuje, netřeba do textu přidávat jakékoliv explicitní formátování, následuje ukázka sazby tohoto a~následujcících odstavců). Je vhodné znovu připomenout, že autor práce by si měl pohlídat, čím končí řádky. Určitě by to neměly být jednohláskové spojky ani předložky. Pro takové případy zde máme tildu (\textasciitilde), kterou \LaTeX vyhodnocuje jako nezlomitelnou mezeru. Příklad použití, ze kterého se můžete poučit, je uveden např. právě na konci tohoto řádku. Konec řádku zde a~další řádek hned za ním, aneb podívejte se do zdrojového kódu za spojku „a“.

Šablona je nastavena na oboustranný tisk. Za tímto účelem obsahuje na svém počátku místy prázdnou stranu, aby např. zadání začínalo na liché stránce (vpravo).


% ============================================================================ %



\part{TEORETICKÁ ČASŤ}
% Kapitola 1: Kontext a motivácia
% Stav: finálna verzia

\section{Kontext a motivácia}

\subsection{Akademický repozitár univerzitnej knižnice}

Inštitucionálny repozitár (akademický repozitár) predstavuje dôležitú súčasť informačnej infraštruktúry vysokej školy. Neplní iba archivačnú funkciu, ale vytvára aj prostredie pre dlhodobé uchovávanie, sprístupňovanie, indexáciu a opakované využívanie vedeckých výstupov a ich metadát. Súčasne slúži ako jeden z bodov prepojenia medzi lokálnymi informačnými systémami univerzity a širšou sieťou otvorenej vedeckej komunikácie. Aktuálne odporúčania pre repozitáre kladú dôraz najmä na kvalitu popisu, interoperabilitu, použitie perzistentných identifikátorov a možnosť strojového sprístupnenia metadát externým službám \cite{COARGoodPractices2022,OpenAIREGuidelines4}.

V prostredí univerzitnej knižnice má inštitucionálny repozitár špecifické postavenie aj preto, že jeho kvalita nezávisí iba od technickej platformy, ale vo veľkej miere od kvality vstupných údajov. Knižnica musí pracovať s bibliografickými záznamami, ktoré vznikli v rôznych externých systémoch, riadia sa odlišnými pravidlami zápisu a len čiastočne zodpovedajú lokálnym požiadavkám inštitúcie. To platí najmä pri evidencii publikačnej činnosti interných autorov univerzity, ktorých výstupy môžu byť publikované mimo interných systémov a mimo vydavateľského prostredia UTB. Z pohľadu tejto práce preto inštitucionálny repozitár nevystupuje len ako cieľové úložisko, ale aj ako kontext, ktorý určuje nároky na presnosť, konzistentnosť a kontrolovateľnosť metadát pred ich uložením.

Zároveň platí, že moderný inštitucionálny repozitár nemožno chápať izolovane. Predpokladá sa, že je súčasťou interoperabilnej siete, v ktorej sa metadáta vystavujú štandardizovaným spôsobom a sú ďalej využívané agregátormi, vyhľadávacími službami a nadväzujúcimi informačnými systémami. Z toho vyplýva, že nekvalitné alebo nejednotné lokálne metadáta nepredstavujú iba interný evidenčný problém, ale negatívne ovplyvňujú aj ďalšie použitie záznamov mimo samotnej knižnice.

\subsection{Súčasný stav workflow knižnice UTB}

Východiskom riešeného problému je aktuálny spôsob spracovania publikačných záznamov v knižnici UTB. Záznamy o vedeckých výstupoch interných autorov sú získavané najmä z externých bibliografických databáz Web of Science a Scopus. Obe platformy poskytujú export záznamov do súborových formátov určených na ďalšie spracovanie, pričom používateľ môže voliť rozsah exportovaných polí a typ výstupného formátu \cite{ClarivateWOSExport,ScopusSupport}. Z technického hľadiska ide o praktický spôsob, ako hromadne získať bibliografické údaje bez potreby manuálneho opisovania jednotlivých záznamov.

Samotný export však ešte nepredstavuje stav vhodný na priamy zápis do internej databázy knižnice. V podmienkach UTB preto nasleduje medzikrok, v ktorom sa záznamy prevedú do tabuľkového formátu a prejdú základnou technickou kontrolou. Následne sa spracúvajú v prostredí Google Sheets, kde pracovníci knižnice ručne opravujú a dopĺňajú polia, ktoré sú neúplné, nekonzistentné alebo nezodpovedajú interným pravidlám evidencie. Až po tejto fáze možno údaje uložiť do internej databázy a ďalej s nimi pracovať v lokálnom prostredí univerzity.

Takto nastavený postup má jednu nespornú výhodu: umožňuje rýchlo využiť existujúce externé zdroje a doplniť ich o odbornú kontrolu knihovníka. Súčasne však vytvára proces, v ktorom sa vedľa seba nachádzajú export z komerčnej databázy, medzivrstva v tabuľkovom nástroji a finálny zápis do interného systému. Čím väčšia je rôznorodosť vstupov, tým viac práce sa presúva na manuálnu interpretáciu a korekciu. Táto situácia sa zhoduje aj s metodickým východiskom Mikuleckého, ktorý pri spracovaní metadát vedeckých výstupov uvažuje o prístupe s jasne oddelenými krokmi extrakcie, validácie a následného spracovania \cite{Mikulecky2025}.

\subsection{Problémové miesta a dôvody automatizácie}

Hlavný problém súčasného workflow nespočíva v samotnej dostupnosti dát, ale v ich nejednotnosti a v náročnosti ich následného zosúladenia s lokálnymi požiadavkami. V exportoch z externých databáz sa opakovane objavujú variantné zápisy mien autorov, neúplné alebo nejednotné afiliácie, skrátené či neštandardizované názvy časopisov, nejednoznačné dátumy redakčného procesu a formálne odchýlky v identifikátoroch. Pri záznamoch bez jednoznačného identifikátora navyše vzniká problém s rozpoznaním duplicít voči už evidovaným položkám. Takéto chyby nemožno odstrániť jedným všeobecným pravidlom, pretože ich povaha je čiastočne formálna a čiastočne sémantická.

Manuálne spracovanie týchto situácií je časovo náročné, repetitívne a zle škáluje s rastúcim objemom záznamov. Záťaž pritom nespočíva iba v samotnom prepise údajov, ale aj v opakovanom rozhodovaní, či dva zápisy označujú tú istú osobu, ten istý časopis alebo tú istú publikáciu. Práve táto interpretačná zložka práce je pre knihovnícke workflow najnáročnejšia. Súčasný proces zároveň len obmedzene podporuje auditovateľnosť, teda spätné určenie toho, ktorý údaj bol zmenený, z akého dôvodu a na základe akého typu zásahu.

Dôvod automatizácie preto nemožno redukovať iba na snahu skrátiť čas spracovania. Rovnako dôležitým cieľom je zvýšenie konzistentnosti výsledných záznamov, zníženie závislosti od ad hoc opráv v tabuľkách a vytvorenie pracovného postupu, v ktorom bude možné odlíšiť technickú validáciu, návrh normalizácie a finálne schválenie človekom. Pre potreby knižnice UTB je pritom rozhodujúce, aby takáto podpora neviedla k nekontrolovanému automatickému prepisu záznamov, ale k lepšie riadenému a overiteľnému workflow.

\subsection{Ciele a štruktúra práce}

Cieľom práce je vytvoriť teoretické a metodické východisko pre návrh webovej aplikácie, ktorá podporí evidenciu a normalizáciu metadát vedeckých publikácií knižnice UTB. Teoretická časť sa preto nesústredí na všeobecný opis repozitárov alebo umelej inteligencie, ale na tie koncepty, ktoré budú mať priamy význam pre praktickú časť. Ide najmä o bibliografické metadáta,  kontrolu autorov, dátovú kvalitu, deduplikáciu záznamov, štandardizáciu dátumov podľa mnormy ISO 8601, auditovateľnosť úprav a o využitie veľkých jazykových modelov ako asistívnej vrstvy pri spracovaní nejednotných vstupov.

Práca zároveň vychádza z dvoch dôležitých lokálnych východísk. Prvým je metodický prístup k spracovaniu a vyhodnocovaniu metadát s podporou AI, ktorý rozpracoval Mikulecký \cite{Mikulecky2025}. Druhým je skúsenosť s návrhom webovej aplikácie nad univerzitnými dátami, vrátane dôrazu na roly používateľov, validáciu a prístupové obmedzenia, ktorú dokumentuje Koňařík \cite{Konarik2025}. Teoretická časť preto pripravuje rámec pre riešenie, ktoré nebude chápať LLM ako autonómneho rozhodovateľa, ale ako podporný komponent integrovaný do kontrolovaného pracovného postupu.

Nasledujúca kapitola sa bude venovať metadátam vedeckých výstupov, autoritnej kontrole a externým zdrojom údajov. Tým sa vytvorí terminologický a konceptuálny základ pre neskoršiu analýzu kvality dát, entity resolution a využitia LLM pri normalizácii záznamov.


% Kapitola 2: Metadáta vedeckých výstupov a authority control
% Stav: finálna verzia

\section{Metadáta vedeckých výstupov a authority control}

Pre návrh aplikácie na evidenciu a normalizáciu publikačných záznamov je potrebné najprv presne vymedziť, aké typy údajov sú predmetom spracovania a akú úlohu plnia v prostredí inštitucionálneho repozitára. Praktická časť nebude pracovať iba s voľným textom, ale s čiastočne štruktúrovanými bibliografickými údajmi rôznej kvality, ktoré sa budú porovnávať, dopĺňať, validovať a prepájať s lokálnymi autoritami. Táto kapitola preto zavádza pojmy a referenčné rámce, bez ktorých by nebolo možné korektne opísať deduplikáciu, normalizáciu ani auditovateľné spracovanie záznamov.

\subsection{Pojem metadát a ich rola v repozitári}

Metadáta možno v najvšeobecnejšom zmysle chápať ako štruktúrované údaje o zdroji. V prostredí vedeckých výstupov ide predovšetkým o údaje, ktoré umožňujú dokument identifikovať, opísať, vyhľadať, previazať s ďalšími entitami a sprístupniť v prepojiteľnom prostredí \cite{DCMITerms}. Pre inštitucionálny repozitár preto metadáta nepredstavujú iba sprievodný opis dokumentu, ale základnú vrstvu, na ktorej stojí evidencia, vyhľadávanie, deduplikácia aj výmena údajov s externými systémami.

V tejto práci sú dôležité najmä tri roviny metadát. Prvú tvoria deskriptívne údaje, ako názov, autori, názov časopisu, rok vydania alebo identifikátory. Druhú rovinu predstavujú väzbové a autoritné údaje, ktoré prepájajú textové polia s konkrétnymi osobami, periodikami alebo inými kontrolovanými entitami \cite{DCMITerms}. Tretiu rovinu tvoria prevádzkové údaje potrebné pre interné spracovanie, napríklad pôvod hodnoty, stav validácie alebo informácia o tom, či bola hodnota doplnená heuristikou, modelom alebo používateľom. Táto tretia rovina je vymedzená pre potreby tejto práce a bude podrobnejšie rozpracovaná v podkapitole venovanej provenancii a auditovateľnosti úprav.

\subsection{Bibliografické metadátové schémy}

Pri spracovaní vedeckých výstupov sa nepoužíva jediná univerzálna schéma. V praxi sa stretávajú všeobecné metadátové modely, knižnične orientované schémy aj prenosové formáty určené na export a výmenu medzi nástrojmi. Pre účely tejto práce je dôležité rozlišovať najmä medzi tým, čo predstavuje konceptuálnu schému opisu, a tým, čo predstavuje technický formát výmeny.

Dublin Core patrí medzi najrozšírenejšie všeobecné metadátové rámce. Jeho význam spočíva najmä v tom, že poskytuje relatívne jednoduchý a prenositeľný súbor prvkov použiteľných naprieč rôznymi typmi repozitárov a informačných služieb \cite{DCMITerms}. Pre bohatší bibliografický opis je však často príliš všeobecný. V tomto smere je podrobnejšia schéma MODS, ktorá vznikla v knižničnom prostredí a umožňuje presnejšie zachytiť bibliografické a vzťahové charakteristiky dokumentu \cite{MODSUserGuide}. Rozdiel medzi týmito prístupmi je pre túto prácu dôležitý preto, že importované záznamy môžu obsahovať údaje vhodné na jednoduchú výmenu, no zároveň vyžadujú detailnejšiu internú reprezentáciu pre potreby kurácie a normalizácie.

Popri metadátových schémach treba rozlišovať aj exportné a citačné formáty. BibTeX je formát pevne spätý s prostredím \LaTeX{} a slúži predovšetkým na uchovávanie a generovanie bibliografických záznamov v citačných workflow \cite{Patashnik1988}. V bibliografických databázach sa však používa aj formát RIS, ktorý v praxi slúži ako rozšírený výmenný formát medzi databázami, citačnými manažérmi a ďalšími nástrojmi. Pre túto prácu nie je rozhodujúca jeho úplná syntaktická špecifikácia, ale skutočnosť, že rôzne vstupné formáty vyjadrujú podobné bibliografické údaje odlišným spôsobom a že aplikácia bude musieť tieto rozdiely preklenúť pri importe a transformácii.

\subsection{Perzistentné identifikátory}

Osobitné miesto v správe metadát majú perzistentné identifikátory. Ich hlavnou úlohou je zabezpečiť, aby bolo možné konkrétnu entitu jednoznačne a dlhodobo identifikovať naprieč systémami, nezávisle od toho, v akej textovej podobe je opísaná v konkrétnom zázname. V rámci riešenej domény ide najmä o identifikátory dokumentov, autorov a seriálových zdrojov.

Pri jednotlivých publikáciách zohráva kľúčovú úlohu DOI. Tento identifikátor je v praxi najspoľahlivejším oporným bodom pri deduplikácii záznamov, pretože je navrhnutý práve na jednoznačnú identifikáciu objektu a jeho strojové rozlíšenie \cite{DOIFoundation,ISO26324_2025}. V prostredí evidencie publikácií preto predstavuje prirodzený primárny kľúč, ak je v zázname dostupný a syntakticky korektný. Pri autoroch má analogickú, hoci nie úplne univerzálnu, úlohu ORCID. Ten nepokrýva všetkých autorov, ale tam, kde je k dispozícii, výrazne znižuje nejednoznačnosť pri párovaní mien a pri prepájaní záznamov medzi systémami \cite{ORCIDWhatIs}. V prípade periodík slúžia ako stabilizačný prvok ISSN a e-ISSN, ktoré umožňujú odlíšiť samotný časopis od variantných textových zápisov jeho názvu \cite{ISSNWhatIs,ISO3297_2022}.

Z pohľadu praktickej časti je dôležité, že identifikátory nevyriešia celú úlohu samy osebe. Nie všetky importované záznamy ich obsahujú, nie vždy sú uvedené v správnom formáte a nie všetky polia sú medzi zdrojmi konzistentné. Napriek tomu práve identifikátory tvoria základný referenčný rámec, voči ktorému sa budú posudzovať variantné zápisy v textových poliach. Ich prítomnosť preto zásadne ovplyvňuje, či sa záznam bude dať riešiť deterministicky, alebo bude potrebné prejsť na porovnanie podobnosti.

\subsection{Authority control}

Authority control predstavuje súbor princípov a postupov, ktorých cieľom je zabezpečiť konzistentné reprezentovanie entít v bibliografickom a repozitárovom prostredí. Nejde iba o technický register mien, ale o spôsob, ako oddeliť samotnú entitu od jej textových variantov a ako nad touto entitou dlhodobo udržiavať kontrolovaný opis. V modeli FRAD je authority data chápané ako základ pre riadené prístupové body a pre vzťahy medzi rôznymi podobami mien, názvov a identít \cite{FRAD2013}.

V kontexte tejto práce sa authority control týka predovšetkým dvoch oblastí. Prvou sú autori, pri ktorých môže mať tá istá osoba viacero zápisov mena, rôzne afiliácie alebo iba čiastočne uvedené iniciály. Druhou sú časopisy, ktorých názvy sa v exportoch môžu objavovať v plnej aj skrátenej podobe, s interpunkčnými odchýlkami alebo s rozdielnym uvedením identifikátorov. Pre praktické riešenie je preto potrebné vychádzať z lokálnych autoritných záznamov a nevnímať textové pole ako definitívny opis entity. Externé služby, ako ORCID, ISSN Portal, Crossref alebo VIAF, môžu slúžiť ako doplnkový referenčný zdroj, no finálna reprezentácia musí byť v lokálnom workflow kontrolovaná vzhľadom na potreby UTB \cite{ORCIDWhatIs,ISSNWhatIs,VIAF,CrossrefRESTAPI}.

Z hľadiska návrhu aplikácie je authority control dôležitý aj preto, že vytvára most medzi normalizáciou textu a evidenčnou integritou systému. Cieľom nie je iba prepísať nejednotný reťazec do esteticky lepšej podoby, ale priradiť ho ku konkrétnej entite, ktorú možno ďalej jednoznačne používať pri vyhľadávaní, validácii, deduplikácii a reportovaní. Táto logika bude v praktickej časti priamo prítomná pri mapovaní autorov na interný register UTB autorov a pri normalizácii časopisov.

\subsection{Externé zdroje pravdy}

Pri správe publikačných záznamov nemožno predpokladať, že jeden zdroj poskytne vždy úplné a bezchybné metadáta. Preto je vhodnejšie hovoriť o externých zdrojoch pravdy v pluráli a chápať ich ako referenčné body s rôznou mierou autority pre rôzne typy údajov. V podmienkach UTB majú osobitné postavenie Web of Science a Scopus, pretože predstavujú hlavné vstupné zdroje importu. Ich výhodou je rozsiahle pokrytie a štandardizované možnosti exportu, nevýhodou však je, že výsledné záznamy nemusia zodpovedať lokálnym požiadavkám na autority, interné väzby ani finálnu podobu evidencie \cite{ClarivateWOSExport,ScopusSupport}.

Významnú doplnkovú rolu má Crossref. Jeho rozhranie REST API umožňuje vyhľadávať a získavať bibliografické metadáta k DOI, periodikám aj jednotlivým dielam vo forme vhodnej na ďalšie strojové spracovanie \cite{CrossrefRESTAPI}. Crossref zároveň explicitne uvádza, že jeho metadáta sú voľne dostupné pre komunitu a že väčšina bibliografických metadát je považovaná za fakty nepodliehajúce autorskému právu, pričom špecifické obmedzenia sa týkajú najmä abstraktov \cite{CrossrefMetadataRetrieval}. Pre riešenú aplikáciu je preto Crossref vhodným doplnkovým zdrojom najmä tam, kde treba overiť alebo doplniť časopis, DOI alebo ďalšie chýbajúce bibliografické polia.

Otvorený charakter má aj OpenAlex, ktorý poskytuje verejne dostupné rozhranie nad rozsiahlym agregovaným grafom vedeckých entít \cite{OpenAlexDocs}. Jeho odborné zarámcovanie poskytuje aj práca Priema, Piwowar a Orra, ktorá OpenAlex opisuje ako plne otvorený index vedeckých diel, autorov, časopisov, inštitúcií a konceptov \cite{OpenAlexPaper}. Z pohľadu tejto práce však OpenAlex nepredstavuje primárny importný kanál, ale skôr pomocný zdroj na overovanie alebo dopĺňanie vybraných údajov. Praktická časť preto nebude vychádzať z predstavy jediného absolútneho zdroja pravdy, ale z kombinácie primárnych importných databáz, otvorených referenčných služieb a lokálnych autoritných záznamov. Práve toto viaczdrojové poňatie je predpokladom pre spoľahlivú normalizáciu metadát v prostredí univerzitnej knižnice.


% Kapitola 3: Kvalita dát a entity resolution
% Stav: finálna verzia

\section{Kvalita dát a entity resolution}

Predchádzajúca kapitola ukázala, že bibliografický záznam nie je iba súborom deskriptívnych polí, ale aj nositeľom identít, väzieb a pravidiel, podľa ktorých sa má ďalej spracovať. Ak má aplikácia podporovať normalizáciu a bezpečný zápis do lokálneho prostredia, nestačí poznať len význam jednotlivých polí. Potrebné je aj určiť, ktoré typy chýb sa v údajoch objavujú, ako možno posudzovať ich kvalitu a akými metódami možno zistiť, že dva zápisy označujú tú istú entitu alebo ten istý záznam. Táto kapitola preto prechádza od opisu metadát k ich analytickému spracovaniu.

\subsection{Rozmery dátovej kvality}

Kvalitu dát nemožno redukovať na jediné kritérium. V literatúre sa opakovane uvádza, že údaje môžu byť formálne správne, no napriek tomu nevhodné na konkrétne použitie. Wang a Strong preto navrhli chápať dátovú kvalitu ako viacrozmerný koncept, ktorého posúdenie závisí od potrieb používateľa a od kontextu použitia \cite{WangStrong1996}. Pre potreby tejto práce sú najrelevantnejšie najmä štyri dimenzie: úplnosť, presnosť, konzistentnosť a časová primeranosť.

Úplnosť vyjadruje, či záznam obsahuje všetky údaje potrebné na zamýšľané použitie. V publikačných exportoch môže byť problémom chýbajúci DOI, neúplný zoznam autorov alebo absentujúca afiliácia. Presnosť sa vzťahuje na mieru zhody medzi údajom v systéme a stavom, ktorý má reprezentovať. Typickým príkladom je nesprávne priradenie časopisu, mylný rok vydania alebo zle rozpoznaná identita autora. Konzistentnosť vyjadruje vnútornú bezrozpornosť údajov a ich zápisu. V bibliografickej oblasti sa prejavuje napríklad rozdielnymi podobami toho istého mena, odlišným formátom identifikátora alebo konfliktom medzi názvom časopisu a ISSN. Časová primeranosť, označovaná aj ako timeliness, vyjadruje, či údaje zodpovedajú aktuálnemu a použiteľnému stavu pre konkrétny pracovný krok \cite{WangStrong1996,PipinoLeeWang2002}.

Z pohľadu navrhovanej aplikácie je dôležité, že tieto dimenzie nemožno posudzovať oddelene. Chýbajúci údaj znižuje úplnosť, no zároveň môže znemožniť presné párovanie a následne znížiť aj konzistentnosť lokálnej evidencie. Abedjan a kol. upozorňujú, že profilovanie a skúmanie datasetu je prirodzeným predpokladom jeho ďalšieho čistenia, integrácie a kontroly \cite{Abedjan2019}. Praktická časť preto nebude chápať kvalitu dát ako abstraktnú vlastnosť databázy, ale ako súbor konkrétnych kontrol a rozhodnutí viazaných na jednotlivé polia publikačného záznamu.

\subsection{Typické chyby v exportoch z bibliografických databáz}

V exportoch z bibliografických databáz sa zvyčajne nekombinuje iba jeden typ chyby, ale viacero vrstiev nejednotnosti naraz. Časté sú chýbajúce hodnoty, variantné textové podoby, konfliktné údaje medzi poliami a technické nedostatky vzniknuté pri prenose alebo transformácii. Takéto chyby majú v bibliografickom kontexte osobitný význam, pretože aj malá odchýlka v mene autora alebo názve periodika môže znemožniť správne priradenie k lokálnej autorite alebo vytvoriť falošnú duplicitu.

Mikulecký pri spracovaní publikačných metadát upozorňuje, že problémové sú najmä polia, ktoré majú v zdrojoch vysokú jazykovú variabilitu, ale v cieľovom systéme musia mať kontrolovanú a štruktúrovanú podobu \cite{Mikulecky2025}. V podmienkach UTB ide najmä o mená autorov, afiliácie, názvy časopisov, redakčné dátumy a identifikátory. Časť problémov je formálna, napríklad nadbytočné medzery, rozdielna interpunkcia alebo nejednotné použitie veľkých písmen. Druhá časť je sémantická, keď dve hodnoty vyzerajú odlišne, ale označujú tú istú entitu, alebo naopak vyzerajú podobne, no vzťahujú sa na odlišné objekty.

Pre návrh aplikácie je dôležité, že tieto chyby nemožno riešiť iba jedným mechanizmom. Niektoré sa dajú odhaliť jednoduchou validáciou formátu, iné vyžadujú porovnanie s externým zdrojom a ďalšie si vyžadujú rozhodovanie nad podobnosťou a kontextom. Práve preto je vhodné oddeliť technickú kontrolu vstupu od úloh entity resolution a od následnej kurácie človekom.

\subsection{Entity resolution}

Pojem entity resolution označuje súbor metód, ktorých cieľom je určiť, či dva alebo viaceré záznamy odkazujú na tú istú reálnu entitu. V literatúre sa používajú aj príbuzné označenia, napríklad record linkage, duplicate detection alebo data matching, no všetky smerujú k rovnakému jadru problému: rozdielne reprezentácie treba previesť na rozhodnutie o identite alebo neidentite \cite{Christen2012,KopckeRahm2010}. V kontexte tejto práce sa entity resolution týka autorov, časopisov aj samotných publikačných záznamov.

\subsubsection{Deterministické pravidlá}

Deterministický prístup vychádza z pravidiel, pri ktorých možno jednoznačne povedať, aký výsledok má porovnanie priniesť. Typickým príkladom je zhodný DOI po normalizácii zápisu alebo presná zhoda s lokálnym identifikátorom entity. Výhodou takéhoto prístupu je jeho predvídateľnosť, auditovateľnosť a nízka interpretačná nejednoznačnosť. Christen uvádza, že presné pravidlá sú vhodné všade tam, kde sú k dispozícii dostatočne spoľahlivé identifikátory alebo dobre definované formátové obmedzenia \cite{Christen2012}.

V bibliografickom workflow však deterministické pravidlá nestačia na všetky situácie. Ich sila spočíva najmä v tom, že dokážu odfiltrovať jednoduché a vysokodôveryhodné prípady. V navrhovanom riešení preto majú tvoriť prvú vrstvu spracovania. Ak je k dispozícii spoľahlivý identifikátor alebo jasné pravidlo, nie je dôvod prechádzať na výpočtovo a interpretačne náročnejšie porovnanie podobnosti.

\subsubsection{Podobnostné metriky}

Ak presné identifikátory chýbajú alebo sú neúplné, nastupuje aproximatívne porovnávanie textových hodnôt. V tejto oblasti sa používajú rôzne triedy podobnostných metrík, ktoré sa líšia citlivosťou na preklepy, transpozície, skrátenia alebo slovosled. Pri menách autorov a názvoch publikácií majú význam najmä edit-distance prístupy a metriky založené na čiastočnej zhode reťazcov. Christen zdôrazňuje, že neexistuje univerzálne najlepšia podobnostná funkcia a jej vhodnosť závisí od typu dát aj od charakteru chyby \cite{Christen2012}.

Pre potreby tejto práce sú relevantné najmä metriky typu Levenshtein a Jaro-Winkler. Editačná vzdialenosť je vhodná tam, kde treba zachytiť malé textové odchýlky vzniknuté vložením, odstránením alebo zámenou znaku \cite{Christen2012}. Jaro-Winklerova podobnosť je zasa praktická pri krátkych reťazcoch, najmä pri menách, kde býva dôležitý spoločný prefix a menšie preklepy v závere slova \cite{Winkler1990}. Z praktického hľadiska však ani podobnostná metrika sama osebe nevyrieši rozhodnutie o zhode. Potrebné je určiť, ktoré polia sa porovnávajú, ako sa ich výsledky kombinujú a pri akom prahu sa už zhoda považuje za dostatočnú.

\subsection{Deduplikácia záznamov}

Deduplikácia predstavuje praktickú aplikáciu entity resolution na úrovni celých záznamov. Jej cieľom nie je len nájsť identické položky, ale zabrániť viacnásobnému uloženiu toho istého publikačného výstupu v dôsledku rozdielnych exportov, odlišného zápisu polí alebo absencie niektorých identifikátorov. Naumann a Herschel upozorňujú, že pri väčších datasetoch je potrebné súčasne riešiť dve otázky: ako posúdiť podobnosť dvoch kandidátnych záznamov a ako obmedziť počet dvojíc, ktoré sa vôbec budú porovnávať \cite{NaumannHerschel2010}.

Prvá otázka súvisí s výberom porovnávacích kľúčov. V podmienkach tejto práce má mať prioritu jednoznačný identifikátor publikácie, teda DOI, ak je dostupný a validný. Ak takýto kľúč chýba, je potrebné použiť kombináciu sekundárnych znakov, najmä názov práce, identifikátor časopisu a rok vydania. Druhá otázka sa týka obmedzenia priestoru kandidátov. Christen aj Naumann a Herschel uvádzajú, že porovnávať všetky záznamy so všetkými je pri väčších súboroch neefektívne, preto sa používajú stratégie blokovania, ktoré najprv vytvoria menšie skupiny pravdepodobných kandidátov \cite{Christen2012,NaumannHerschel2010}.

V navrhovanom riešení bude mať deduplikácia význam nielen voči už existujúcim záznamom repozitára, ale aj medzi importmi z rôznych zdrojov. Praktický problém preto nie je binárny, ale viacvrstvový: treba rozlíšiť zhodu v rámci jedného importu, medzi dvoma importnými dávkami a medzi novou dávkou a cieľovým úložiskom. Teoretickým východiskom je pritom kombinácia presných pravidiel a podobnostného porovnávania, nie voľné rozhodovanie bez auditovateľného základu.

\subsection{Štandardizácia časových údajov}

Osobitný problém predstavujú časové údaje, ktoré sa v bibliografických záznamoch často vyskytujú vo voľnom texte alebo v nejednotnej štruktúre. V publikačnom workflow môžu mať význam nielen klasické údaje o roku vydania, ale aj dátumy redakčného procesu, napríklad prijatie rukopisu, revízia, akceptovanie alebo prvé zverejnenie. Ak sa majú tieto hodnoty ďalej validovať, porovnávať alebo ukladať do databázy, je potrebné previesť ich do jednotného a strojovo jednoznačného formátu.

Základným referenčným rámcom je norma ISO 8601-1:2019, ktorá určuje pravidlá reprezentácie dátumu a času pre informačnú výmenu \cite{ISO8601_2019}. V aplikačnej praxi je veľmi relevantný aj RFC 3339, ktorý predstavuje internetový profil týchto zápisov a presne určuje formu timestampu vrátane časovej zóny \cite{RFC3339}. Pre potreby tejto práce je dôležité najmä to, že jednotný formát odstraňuje nejednoznačnosť pri ďalšom spracovaní. Ak sa dátumy ponechajú iba vo voľnom texte, stráca sa možnosť ich spoľahlivo porovnávať, filtrovať a validovať.

Štandardizácia časových údajov preto nie je iba estetická úprava formátu. Ide o prevod z ľudsky čitateľnej, ale pre systém často nejednoznačnej reprezentácie do tvaru, ktorý možno jednoznačne uložiť, kontrolovať a odovzdať ďalším vrstvám aplikácie. Tento krok bude mať v praktickej časti osobitné postavenie, pretože kombinuje parsovanie textu, doménové pravidlá a následnú normalizáciu do jednotného zápisu.

\subsection{Provenancia a auditovateľnosť úprav}

Ak má aplikácia podporovať zásahy do importovaných metadát, nestačí ukladať iba finálnu hodnotu po oprave. Dôležité je aj to, aby bolo možné spätne určiť, odkiaľ daná hodnota pochádza, akým spôsobom bola získaná a kto alebo čo ju navrhlo. W3C PROV definuje provenanciu ako informáciu o entitách, aktivitách a osobách zapojených do vzniku alebo zmeny údajov, pričom práve takáto informácia umožňuje posudzovať kvalitu, spoľahlivosť a dôveryhodnosť dát \cite{PROVOverview}.

V doméne správy publikačných metadát má provenancia priamy praktický význam. Pri jednom poli môže existovať pôvodná hodnota z importu, návrh doplnenia z externého zdroja, odporúčanie generované modelom a napokon manuálne potvrdená alebo upravená hodnota knihovníkom. Bez evidencie pôvodu takýchto zásahov nie je možné spoľahlivo analyzovať, ktoré mechanizmy fungujú dobre, kde vznikajú chyby a aký bol dôvod výslednej podoby záznamu. Mikulecký v metodickom rámci práce s AI pri metadátach zdôrazňuje potrebu sledovať úspešnosť a interpretovateľnosť jednotlivých krokov spracovania \cite{Mikulecky2025}. V podobnom duchu Koňařík pri webovej aplikácii nad univerzitnými dátami akcentuje validáciu a kontrolovaný používateľský zásah ako súčasť dôveryhodného workflow \cite{Konarik2025}.

Pre praktickú časť z toho vyplýva, že auditovateľnosť nemožno chápať iba ako technický log operácií. Má ísť o vedomé modelovanie pôvodu hodnôt a typov zásahu tak, aby bolo možné odlíšiť pôvodný import, heuristické rozhodnutie, odporúčanie AI a finálne schválenie človekom. Takáto evidencia je dôležitá nielen pre spätnú kontrolu, ale aj pre vyhodnocovanie kvality navrhovaného riešenia.


% Kapitola 4: LLM ako podpora správy metadát
% Stav: finálna verzia

\section{LLM ako podpora správy metadát}

Predchádzajúca kapitola ukázala, že pri správe publikačných metadát sa opakovane objavujú situácie, ktoré nemožno uspokojivo vyriešiť len jednoduchou syntaktickou validáciou. Ide najmä o prípady, kde treba interpretovať voľný text, porovnávať variantné zápisy, rozhodovať medzi viacerými kandidátmi alebo prevádzať nejednotné vstupy do kontrolovanej štruktúry. Práve v tomto priestore sa otvára miesto pre využitie veľkých jazykových modelov ako asistenčnej vrstvy nad tradičnými heuristikami a validačnými pravidlami.

Táto kapitola preto nepojednáva o LLM ako o všeobecnom fenoméne, ale výlučne v rozsahu potrebnom pre túto prácu. Dôraz sa kladie na tie vlastnosti modelov, ktoré sú priamo relevantné pre extrakciu, normalizáciu a kontrolu bibliografických údajov, na návrh ich vstupov a výstupov a na podmienky, za ktorých zostáva ich nasadenie bezpečné, auditovateľné a metodicky obhájiteľné.

\subsection{Stručný princíp LLM}

Veľké jazykové modely sú generatívne modely trénované na rozsiahlych textových korpusoch, ktoré dokážu na základe vstupného kontextu predikovať a vytvárať ďalší text. Pre potreby tejto práce je dôležité najmä to, že model dokáže pracovať s neštruktúrovaným alebo len čiastočne štruktúrovaným vstupom a prevádzať ho do podoby, ktorá lepšie zodpovedá požiadavkám následného strojového spracovania \cite{AlammarGrootendorst2024,Foster2019}. V praxi to znamená, že model môže pomôcť tam, kde sa v jednom poli mieša faktická informácia, prirodzený jazyk a nejednotný zápis.

Zároveň však ide o model pravdepodobnostný, nie deterministický. To znamená, že ani pri rovnakom zadaní nemožno bez ďalších opatrení automaticky predpokladať úplne totožný alebo bezchybný výstup. Yanev zdôrazňuje, že pri budovaní aplikácií nad LLM je potrebné navrhovať ich použitie tak, aby sa ich výstupy dali validovať, obmedziť a spoľahlivo začleniť do širšej aplikačnej logiky \cite{Yanev2024}. V kontexte tejto práce je preto podstatné chápať LLM nie ako náhradu evidenčného systému, ale ako komponent určený na podporu konkrétnych kurátorských úloh.

\subsection{Roly LLM v metadátovom workflow}

V metadátovom workflow možno jazykový model chápať minimálne v troch odlišných rolách. Prvou je rola normalizátora, teda nástroja, ktorý z nejednotného textového vstupu vytvára kandidátnu normalizovanú hodnotu alebo štruktúrovaný návrh zodpovedajúci požiadavkám systému. Druhou je rola asistenta rozhodovania, v ktorej model pomáha používateľovi pri výbere medzi viacerými kandidátmi alebo pri interpretácii nejednoznačného poľa. Tretia rola je rola sekundárneho kontrolóra konzistencie, keď model upozorňuje na možné rozpory medzi poliami alebo na neúplnosť záznamu \cite{Mikulecky2025}.

Tieto roly sa navzájom odlišujú mierou rizika aj typom očakávaného výstupu. Pri normalizácii je cieľom najmä transformácia voľného textu do kontrolovanej podoby. Pri asistencii rozhodovania je dôležitejšia argumentovaná voľba medzi alternatívami. Pri kontrole konzistencie ide zasa o schopnosť modelu zachytiť, že niektoré polia si navzájom odporujú alebo že záznam nezodpovedá očakávanej logike. Mikulecký ukazuje, že práve úlohy viazané na extrakciu a transformáciu polí patria medzi najsľubnejšie oblasti nasadenia AI v evidencii vedeckej publikačnej činnosti \cite{Mikulecky2025}. Pre túto prácu je však rozhodujúce, že vo všetkých troch rolách musí model fungovať ako asistent v rámci kontrolovaného workflow, nie ako autonómny vykonávateľ zápisov do databázy.

\subsection{Štruktúrované výstupy a prompt engineering}

Pri spracovaní bibliografických metadát nestačí, aby model vytvoril jazykovo presvedčivú odpoveď. Výstup musí byť navrhnutý tak, aby sa dal spoľahlivo parsovať, validovať a uložiť do aplikácie. Preto je v tejto doméne dôležitejšia schopnosť generovať štruktúrované odpovede než schopnosť formulovať voľný text. OpenAI aj Ollama dnes explicitne podporujú structured outputs, teda režim, v ktorom sa odpoveď modelu podriaďuje zadanej JSON schéme \cite{OpenAIStructuredOutputs,OllamaStructuredOutputs}. Takýto prístup významne znižuje rozdiel medzi jazykovým modelom a bežnou aplikačnou logikou, pretože výstup modelu možno bezprostredne podrobiť formálnej validácii.

S tým úzko súvisí aj prompt engineering, teda návrh vstupov tak, aby bol model vedený k čo najstabilnejšiemu a najrelevantnejšiemu správaniu. V tejto práci však nejde o prompt engineering v širokom zmysle, ale o jeho úzko praktickú podobu zameranú na extrakciu polí, dopĺňanie kandidátnych hodnôt a kontrolu konzistencie. Yanev aj Alammar s Grootendorstom sa zhodujú, že pri aplikáciách orientovaných na dáta je potrebné zadanie formulovať presne, minimalizovať nejednoznačnosť a explicitne špecifikovať požadovaný výstupný formát \cite{Yanev2024,AlammarGrootendorst2024}.

\subsubsection{JSON schéma a function calling}

Štruktúrovaný výstup založený na JSON schéme má v tejto práci zásadný význam. Umožňuje totiž presne určiť, aké polia sa od modelu očakávajú, aký majú mať typ, ktoré z nich sú povinné a aké obmedzenia sa na ne vzťahujú. OpenAI rozlišuje structured outputs pre odpovede priamo v JSON schéme a function calling, pri ktorom model vytvára argumenty pre vopred definovaný nástroj alebo funkciu \cite{OpenAIStructuredOutputs,OpenAIFunctionCalling}. Ollama podporuje obdobný princíp tým, že modelu možno pri chat rozhraní odovzdať JSON schému cez parameter \texttt{format} \cite{OllamaStructuredOutputs}.

Pre navrhované riešenie je tento prístup dôležitý preto, že oddeľuje jazykovú interpretáciu vstupu od následnej systémovej akcie. Model môže napríklad navrhnúť štruktúrovaný objekt s menom autora, kandidátnou afiliáciou alebo normalizovaným dátumom, no až ďalšia vrstva aplikácie rozhodne, či je tento návrh formálne validný, či je v zhode s autoritami a či bude predložený používateľovi na schválenie. Function calling je v tomto kontexte užitočný najmä preto, že model nevracia iba text, ale argumenty pre presne vymedzenú operáciu, čo zvyšuje kontrolovateľnosť celého procesu.

\subsubsection{Prompt engineering}

Prompt engineering má v tejto doméne zmysel predovšetkým ako nástroj obmedzenia variability modelových výstupov. Oficiálna dokumentácia OpenAI zdôrazňuje potrebu klásť inštrukcie jednoznačne, jasne oddeľovať inštrukcie od vstupných dát, byť konkrétny v požadovanom výstupe a pri zložitejších úlohách používať aj few-shot príklady \cite{OpenAIPromptGuide}. Pri metadátových úlohách je dôležité, aby model presne vedel, ktoré časti vstupu má chápať ako dáta na analýzu a ktoré ako pravidlá spracovania.

Z praktického hľadiska to znamená, že systémový prompt má vymedziť úlohu modelu a hranice jeho správania, zatiaľ čo používateľský vstup alebo importovaný text predstavuje materiál na spracovanie. Few-shot príklady sú užitočné najmä pri poliach s opakujúcou sa štruktúrou, napríklad pri extrakcii dátumov z redakčných poznámok alebo pri mapovaní autorov na očakávanú výstupnú štruktúru. Cieľom nie je vytvoriť univerzálny prompt pre všetky situácie, ale navrhnúť menší počet stabilných promptových šablón pre konkrétne typy úloh.

\subsection{Human-in-the-loop a riziká}

Použitie LLM pri správe metadát má zmysel iba vtedy, ak zostáva človek súčasťou rozhodovacieho procesu. Human-in-the-loop v tomto kontexte neznamená iba formálnu možnosť zásahu, ale vedomý návrhový princíp, podľa ktorého model pripravuje návrh a používateľ nesie zodpovednosť za jeho potvrdenie alebo odmietnutie. Tento princíp je dôležitý najmä tam, kde model pracuje s nejednoznačnými údajmi, s variantnými zápismi autorov alebo s dopĺňaním hodnôt, ktoré nemožno deterministicky overiť \cite{Mikulecky2025}.

Dôvodom nie je len metodická opatrnosť, ale aj bezpečnostné riziká. OWASP upozorňuje, že LLM aplikácie sú zraniteľné voči prompt injection, pri ktorej útočník alebo škodlivý vstup manipuluje správanie modelu prostredníctvom prirodzeného jazyka \cite{OWASPPromptInjection}. Zároveň ide o systémy náchylné na halucinácie, teda na produkciu presvedčivo znejúcich, no vecne nesprávnych hodnôt. V širšom bezpečnostnom rámci OWASP patrí medzi relevantné riziká aj únik citlivých informácií, neprimeraná dôvera v modelové výstupy, nedostatočná kontrola nástrojov a závislosť od dodávateľského reťazca modelov a knižníc \cite{OWASPTop10LLM2025}. V doméne bibliografických metadát síce zvyčajne nejde o vysoko citlivé osobné údaje, no aj nesprávne alebo nekontrolovane prijaté výstupy môžu poškodiť kvalitu evidencie a dôveryhodnosť systému.

Preto sa ako vhodný model javí taký návrh, v ktorom LLM navrhuje, no nekoná sám. Výstup modelu má byť validovaný, auditovaný a v rizikových prípadoch predkladaný používateľovi ako kandidátna hodnota. Takýto prístup znižuje nielen pravdepodobnosť chybného zápisu, ale aj prevádzkové riziká spojené s tým, že model by mal priamy a nekontrolovaný dosah na interné údaje.

\subsection{Lokálne vs cloud nasadenie}

Pri integrácii LLM do informačného systému treba rozhodnúť aj o spôsobe nasadenia. V zásade ide o voľbu medzi lokálnym režimom, v ktorom model beží v prostredí organizácie alebo vývojára, a cloudovým režimom, v ktorom sa model používa prostredníctvom vzdialeného API. Ollama reprezentuje typický príklad lokálneho prístupu, pri ktorom je model obsluhovaný cez lokálne rozhranie a možno ho začleniť do aplikácie bez závislosti od vzdialeného poskytovateľa \cite{OllamaDocs}. Na opačnej strane stoja cloudové služby s OpenAI-kompatibilným rozhraním, pri ktorých je výhodou jednoduchšia integrácia a často aj vyšší výkon alebo širší výber modelov; príkladom je dokumentovaná kompatibilita Groq API s klientskymi knižnicami OpenAI \cite{GroqOpenAICompatibility}.

Rozhodovanie medzi týmito režimami nemožno zužovať na jediný parameter. Lokálne nasadenie znižuje závislosť od externého poskytovateľa a zjednodušuje kontrolu nad tokom dát. Cloudové nasadenie zasa často ponúka nižšie bariéry pri škálovaní, menšiu prevádzkovú záťaž a prístup k modelom, ktoré by lokálne nebolo praktické prevádzkovať. Yanev pri budovaní aplikácií nad LLM upozorňuje, že architektúra systému musí zohľadniť nielen kvalitu modelu, ale aj latenciu, spoľahlivosť rozhrania, cenu používania a spôsob validácie odpovedí \cite{Yanev2024}. V tejto práci preto nejde o normatívny výber jednej možnosti, ale o návrhový rámec, v ktorom môže byť LLM vrstva vymeniteľná podľa prevádzkových potrieb.

\subsection{Vyhodnocovanie úspešnosti}

Ak má byť použitie LLM v práci metodicky obhájiteľné, nestačí konštatovať, že model „pomáha“. Potrebné je určiť, v ktorých úlohách sa jeho prínos meria, aký je referenčný stav a akým spôsobom sa bude vyhodnocovať správnosť alebo užitočnosť výstupu. V kontexte tejto práce je prirodzené hodnotiť úspešnosť na úrovni jednotlivých polí, nie iba na úrovni celého záznamu. Osobitne možno posudzovať napríklad správnosť identifikácie interného autora, normalizácie názvu časopisu, extrakcie dátumu alebo návrhu deduplikačného rozhodnutia.

Mikulecký pri hodnotení AI podpory evidencie publikačných výstupov pracuje s porovnaním voči očakávanému výsledku a ukazuje, že relevantné je sledovať nielen všeobecnú úspešnosť, ale aj typy chýb viazané na konkrétne polia \cite{Mikulecky2025}. V oblasti párovania a entity resolution sa tradične používajú metriky typu precision a recall, prípadne ich kombinácie, pričom Christen zdôrazňuje význam kvalitne pripraveného referenčného súboru, voči ktorému sa výsledok porovnáva \cite{Christen2012}. Pre túto prácu z toho vyplýva potreba gold standardu, teda manuálne overenej vzorky záznamov, na ktorej bude možné overiť prínos navrhovaného riešenia.

Pri úlohách, ktoré obsahujú určitú mieru interpretačnej nejednoznačnosti, je vhodné sledovať aj zhodu medzi ľudskými anotátormi. Cohenova kappa je klasická miera zhody pre nominálne kategórie a umožňuje odlíšiť skutočnú zhodu od zhody, ktorá by mohla vzniknúť náhodne \cite{Cohen1960}. V kontexte tejto práce má takáto miera význam najmä pri posudzovaní toho, či je určitý záznam duplicitný, či daný autor patrí medzi interných autorov alebo či je navrhnutá normalizácia ešte prijateľná z pohľadu knihovníckej praxe. Vyhodnocovanie úspešnosti preto nebude iba technickým doplnkom implementácie, ale nevyhnutnou súčasťou argumentácie, že integrácia LLM do workflow prináša merateľný prínos.


% Kapitola 5: Súvisiace systémy a porovnateľné riešenia
% Stav: finálna verzia

\section{Súvisiace systémy a porovnateľné riešenia}

Pri návrhu aplikácie na podporu evidencie a normalizácie publikačných metadát je užitočné zasadiť riešený problém do širšieho prostredia existujúcich systémov. Cieľom tejto kapitoly nie je podávať úplný prehľad všetkých repozitárových a výskumno-informačných platforiem, ale vybrať tie typy systémov, ktoré sú pre tému práce porovnateľné z hľadiska funkcie, dátového modelu alebo pracovného postupu. Rozlíšenie medzi repozitárom, CRIS systémom a otvoreným agregátorom je dôležité aj preto, že navrhovaná aplikácia sa nachádza medzi týmito kategóriami iba čiastočne a nemožno ju bez ďalšieho stotožniť so žiadnou z nich.

\subsection{Inštitucionálne repozitáre}

Najbližším referenčným typom systému je inštitucionálny repozitár. Jeho jadrom je uchovávanie, sprístupňovanie a dlhodobá správa digitálnych objektov a ich metadát v prostredí konkrétnej organizácie. Z praktického hľadiska ide o platformu, ktorá spravidla kombinuje evidenciu záznamov, ukladanie súborov, prístupové politiky, export metadát a nadväzujúce integračné mechanizmy. Pre túto prácu sú relevantné najmä DSpace, EPrints a InvenioRDM, pretože reprezentujú tri etablované open-source prístupy k repozitárovej infraštruktúre.

DSpace patrí medzi najznámejšie a najrozšírenejšie repozitárové riešenia. Už jeho pôvodný opis zdôrazňoval úlohu digitálneho repozitára ako inštitucionálnej služby na uchovávanie a redistribúciu výskumných a vzdelávacích materiálov \cite{Smith2003}. Súčasná dokumentácia DSpace ukazuje, že platforma je orientovaná na správu komunít, kolekcií, metadát, autorizácie, exportu a ingestu obsahu a že podporuje integrácie s externými službami vrátane ORCID a OpenAIRE \cite{DSpaceDocs}. Pre potreby tejto práce je dôležité, že DSpace reprezentuje model plnohodnotného repozitára, nie iba pomocného validačného nástroja pred uložením záznamu.

EPrints predstavuje iný, historicky veľmi významný open-source prístup. Oficiálna dokumentácia aj produktová prezentácia zdôrazňujú flexibilitu workflow, prispôsobiteľnosť formulárov, lokálnu konfiguráciu schémy a použitie v prostredí otvorených repozitárov aj iných typov inštitucionálneho zberu obsahu \cite{EPrintsWhatIs}. Z oficiálnej dokumentácie EPrints sa zdá, že platforma kladie výraznejší dôraz na akvizičný a kurátorský workflow \cite{EPrintsWhatIs}, čo je z hľadiska tejto práce zaujímavé, pretože riešená aplikácia pracuje práve v štádiu pred finálnym uložením záznamu.

InvenioRDM reprezentuje novší typ repozitárovej platformy, ktorá je od začiatku koncipovaná ako turn-key research data management repository s dôrazom na API, interoperabilitu, škálovateľnosť a konfigurovateľnosť \cite{InvenioRDMDocs}. Oficiálna dokumentácia uvádza, že InvenioRDM je harvestovateľný cez OAI-PMH, vystavuje REST API, podporuje viacero exportných formátov a interný metadátový model opiera o schému DataCite \cite{InvenioRDMInteroperable}. Z hľadiska tejto práce je cenné najmä to, že ukazuje spojenie repozitárovej logiky s modernou aplikačnou architektúrou. Ani InvenioRDM však nerieši presne ten istý problém ako navrhovaná aplikácia, pretože jeho primárnou úlohou je správa a publikovanie záznamov a objektov, nie kurátorské spracovanie importov pred zápisom do cudzieho lokálneho systému.

\subsection{Komerčné CRIS systémy}

Od inštitucionálnych repozitárov treba odlíšiť CRIS alebo RIMS systémy, teda platformy orientované na zhromažďovanie, prepájanie a reportovanie výskumných informácií na úrovni organizácie. Zatiaľ čo repozitár sa typicky sústreďuje na záznam a sprístupnenie konkrétnych objektov, CRIS systém má širší záber a prepája publikácie, osoby, projekty, granty, organizačné jednotky, hodnotenie a verejné profily.

Typickým príkladom je Pure, ktorý Elsevier charakterizuje ako Research Information Management System alebo Current Research Information System navrhnutý na integráciu rôznych zdrojov výskumných informácií, reporting a prezentáciu výskumných aktivít inštitúcie \cite{PureRIMS}. Podobný profil má Symplectic Elements, ktorý podľa oficiálneho opisu ingestuje dáta z viacerých interných a externých zdrojov a vytvára pre organizáciu jednotný obraz výskumnej činnosti, použiteľný pre profily, reporting a administratívne procesy \cite{SymplecticElements}. Obe platformy sú pre túto prácu zaujímavé skôr ako referenčný kontrast než ako priama technologická inšpirácia.

Podstatné je, že CRIS systémy riešia širší inštitucionálny problém než ten, ktorý je predmetom tejto diplomovej práce. Navrhovaná aplikácia sa nesnaží centralizovať všetky výskumné informácie univerzity, ale zamerať sa na konkrétny medzikrok medzi importom bibliografických záznamov a ich finálnym uložením do existujúceho lokálneho prostredia. V tomto zmysle je vhodné vnímať CRIS systémy ako konceptuálne príbuzné, no rozsahom a ambíciou odlišné riešenia.

\subsection{Otvorené agregátory metadát}

Tretiu skupinu predstavujú otvorené agregátory metadát. Ich úlohou nie je primárne prevádzkovať lokálny repozitár ani slúžiť ako interný organizačný systém, ale zbierať, prepájať a vystavovať metadáta z veľkého počtu zdrojov. Pre túto prácu sú relevantné najmä OpenAIRE a OpenAlex.

OpenAIRE vystupuje ako európska infraštruktúra založená na interoperabilnom zbere metadát z repozitárov, dátových archívov a CRIS systémov. Jeho oficiálne guidelines definujú, ako majú lokálne systémy vystavovať metadáta cez OAI-PMH tak, aby ich bolo možné agregovať, validovať a začleniť do širšej infraštruktúry \cite{OpenAIREGuidelines4}. Z pohľadu tejto práce má OpenAIRE význam najmä ako normatívny rámec pre výmenu a zber metadát, nie ako nástroj na lokálnu kuráciu importovaných záznamov.

OpenAlex má odlišný charakter. Ide o plne otvorený index vedeckých diel, autorov, zdrojov, inštitúcií a konceptov, sprístupnený cez verejné API a dátové snapshoty \cite{OpenAlexPaper,OpenAlexDocs}. Jeho význam pre túto prácu spočíva v tom, že môže slúžiť ako doplnkový externý zdroj na overovanie alebo dopĺňanie vybraných bibliografických polí. Na rozdiel od OpenAIRE však nejde o infraštruktúru založenú na validačných odporúčaniach pre lokálne repozitáre, ale o otvorený metadátový index nad vedeckým ekosystémom.

\subsection{Prečo UTB workflow dopĺňa, nenahrádza repozitár}

Porovnanie s repozitármi, CRIS systémami a agregátormi ukazuje, že navrhované riešenie nemožno presne zaradiť ani do jednej z týchto kategórií. Nejde o plnohodnotný repozitár, pretože jeho cieľom nie je dlhodobé uchovávanie digitálnych objektov, verejné sprístupňovanie záznamov ani implementácia celej publikačnej infraštruktúry. Nejde ani o CRIS systém, pretože sa nesústredí na úplnú správu výskumných informácií organizácie. Rovnako nejde o agregátor, keďže jeho úlohou nie je zber metadát z veľkého počtu externých zdrojov do jednej otvorenej nadstavby.

Navrhovaná aplikácia má užší, ale v lokálnom prostredí dôležitý účel. Predstavuje kurátorskú a validačnú medzivrstvu medzi importom z externých bibliografických databáz a finálnym uložením záznamu do interného prostredia knižnice UTB. Práve v tomto štádiu vzniká potreba normalizácie autorov, časopisov, dátumov, deduplikácie a zaznamenávania provenancie úprav. Repozitár ani CRIS systém tieto úlohy môžu čiastočne obsahovať \cite{DSpaceDocs,PureRIMS}, no nie vždy ich riešia spôsobom zodpovedajúcim lokálnemu workflow konkrétnej inštitúcie.

Z tohto dôvodu je presnejšie hovoriť o doplnkovej aplikačnej vrstve než o náhrade existujúceho repozitára. Praktická časť práce preto nebude smerovať k vytvoreniu nového repozitárového systému, ale k návrhu špecializovaného nástroja, ktorý zlepší kvalitu vstupných záznamov ešte pred ich uložením do existujúcej evidenčnej a repozitárovej infraštruktúry UTB.


% Kapitola 6: Právne a etické aspekty
% Stav: finálna verzia

\section{Právne a etické aspekty}

Právne a etické aspekty navrhovaného riešenia nemožno oddeliť od jeho technického návrhu. Aj keď aplikácia pracuje prevažne s bibliografickými metadátami a nie s plnými textami alebo s osobitnými kategóriami osobných údajov, stále ide o údaje, ktoré môžu byť viazané na identifikované alebo identifikovateľné fyzické osoby. Zároveň sa pri importe a ďalšom spracovaní pracuje s metadátami pochádzajúcimi z viacerých externých zdrojov, pri ktorých sa líšia podmienky opätovného použitia. Táto kapitola preto stručne vymedzuje dve oblasti: spracúvanie údajov o autoroch podľa GDPR a licenčný režim bibliografických metadát \cite{GDPR2016,CrossrefMetadataRetrieval}.

\subsection{Spracúvanie údajov o autoroch a GDPR}

Nariadenie (EÚ) 2016/679 definuje osobné údaje ako akékoľvek informácie týkajúce sa identifikovanej alebo identifikovateľnej fyzickej osoby a spracúvanie ako akúkoľvek operáciu alebo súbor operácií vykonávaných s osobnými údajmi, či už automatizovanými alebo neautomatizovanými prostriedkami \cite{GDPR2016}. V prostredí bibliografických záznamov to znamená, že meno autora, jeho identifikátor ORCID, afiliácia alebo kombinácia viacerých údajov umožňujúcich jednoznačné priradenie ku konkrétnej osobe môžu mať povahu osobných údajov. Samotná práca s publikačnými metadátami preto nestojí mimo rámca ochrany osobných údajov len preto, že ide o údaje vedeckej komunikácie.

Pre navrhovanú aplikáciu je relevantné najmä to, že GDPR nestanovuje len definície, ale aj základné zásady spracúvania osobných údajov. Medzi ne patrí zákonnosť, korektnosť a transparentnosť, účelové viazanie, minimalizácia údajov, presnosť a primeraná bezpečnosť spracovania \cite{GDPR2016}. Tieto zásady dopĺňa aj výklad EDPB k článku 25, ktorý zdôrazňuje ochranu údajov už pri návrhu systému a v predvolených nastaveniach \cite{EDPBGuidelinesDPbDD2019}. V praktickom preklade pre túto prácu to znamená, že aplikácia by mala spracúvať iba tie údaje o autoroch, ktoré sú potrebné pre evidenciu a normalizáciu publikačných záznamov, mala by umožniť opravu nepresných údajov a mala by podporovať kontrolovaný prístup k rozhraniu a auditovateľnosť zásahov. Keďže predmetom spracovania nie sú súkromné texty autorov, ale najmä verejne dostupné bibliografické údaje, jadrom problému nie je utajenie obsahu, ale korektné a primerané nakladanie s identifikačnými údajmi a s históriou ich úprav.

Z etického hľadiska je dôležité aj to, aby automatizovaná podpora normalizácie neviedla k neodôvodnenému prepájaniu alebo zámene autorov. Chybná identifikácia interného autora, nesprávne priradenie afiliácie alebo nekritické prevzatie návrhu modelu môže mať dopad na evidenciu publikačnej činnosti a na jej ďalšie využitie v interných procesoch univerzity. Preto má význam, aby aj v prípadoch, keď sú vstupné údaje verejne dostupné, zostal zachovaný princíp ľudského schvaľovania a spätnej dohľadateľnosti rozhodnutí.

\subsection{Licencovanie a použiteľnosť bibliografických metadát}

Pri práci s externými zdrojmi je potrebné odlíšiť technickú dostupnosť metadát od právnych podmienok ich ďalšieho použitia. V prípade Crossref je situácia pomerne priaznivá. Oficiálna dokumentácia uvádza, že metadáta sprístupňované cez Crossref sú voľne dostupné komunite, že väčšina bibliografických metadát je považovaná za fakty nepodliehajúce autorskému právu a že Crossref-generované údaje sprístupňuje ako public domain material, pričom špecifická pozornosť sa venuje najmä abstraktom, ktoré môžu zostať predmetom autorskoprávnej ochrany vydavateľa alebo autora \cite{CrossrefMetadataRetrieval}. Pre navrhovanú aplikáciu z toho vyplýva, že použitie Crossref ako doplnkového zdroja na overovanie a dopĺňanie bibliografických polí je z licenčného hľadiska podstatne menej problematické než pri uzavretých komerčných databázach.

Odlišná situácia platí pri platformách Web of Science a Scopus. Hoci tieto služby poskytujú používateľom možnosť exportovať bibliografické záznamy, podmienky ich ďalšieho použitia sa odvíjajú od licenčných a zmluvných pravidiel konkrétnej služby a inštitucionálneho prístupu. Clarivate viaže používanie svojich produktov a služieb na všeobecné podmienky a ďalšie produktovo špecifické podmienky \cite{ClarivateTermsOfUse}. Elsevier pri vývojárskom prístupe výslovne uvádza, že jednotlivé use cases a politiky určujú, ako možno dáta získané cez jeho API používať, a že ďalšie oprávnenia závisia od typu používateľa a predplatného príslušnej služby \cite{ElsevierAPIPolicy}. Z pohľadu tejto práce je preto vhodné vychádzať z konzervatívneho predpokladu, že exportované dáta z týchto zdrojov sa majú používať v rozsahu zodpovedajúcom oprávnenému internému workflow univerzity a nie ako voľne redistribuovateľný otvorený dataset.

Rozlíšenie medzi otvorenými a licenčne viazanými zdrojmi má aj etický rozmer. Ak má aplikácia prepájať viacero zdrojov a dopĺňať chýbajúce hodnoty, musí byť zrejmé, z ktorého zdroja údaj pochádza a za akých podmienok bol získaný. Evidencia provenancie preto neslúži iba na technickú auditovateľnosť, ale aj na transparentné odlíšenie údajov prevzatých z otvorene použiteľných služieb od údajov pochádzajúcich z licenčne obmedzeného prostredia. V tomto zmysle právne a etické požiadavky priamo podporujú návrhový princíp, podľa ktorého aplikácia nemá nekriticky zlúčiť všetky dostupné údaje do jednej neodlišenej vrstvy, ale má zachovať ich pôvod, kontext a režim použitia \cite{CrossrefMetadataRetrieval}.



\part{PRAKTICKÁ ČASŤ}
% Kapitola 7: Analýza zdrojových dát a požiadavky na aplikáciu
% Stav: draft

\section{Analýza zdrojových dát a požiadavky na aplikáciu}

Praktická časť práce nadväzuje na teoretické vymedzenie problému z predchádzajúcich kapitol a smeruje k návrhu konkrétneho riešenia pre knižnicu UTB. Aby tento návrh nebol odtrhnutý od reálnych obmedzení, je potrebné najprv presne opísať, s akými dátami a v akom prostredí bude aplikácia pracovať. Táto kapitola preto identifikuje vstupné zdroje, analyzuje súčasný spôsob práce knihovníkov a formuluje funkčné a nefunkčné požiadavky, voči ktorým bude možné neskôr vyhodnotiť úspešnosť navrhnutého riešenia. Súčasťou je aj vedomé vymedzenie hraníc projektu, teda toho, čo aplikácia podporovať bude a čo zámerne nepokrýva.

\subsection{Vstupné zdroje dát}

Aplikácia spracúva metadáta publikačných záznamov, ktoré pochádzajú z viacerých systémov s odlišnou povahou, štruktúrou a kvalitou. Ich rozdelenie je dôležité, pretože v ďalších kapitolách bude každý zdroj plniť odlišnú rolu pri normalizácii, validácii a rozhodovaní o identite záznamov. Tabuľka~\ref{tab:vstupne-zdroje} sumarizuje jednotlivé zdroje, ich formát, spôsob prístupu a rolu, ktorú v workflow plnia. Podrobnejší popis nasleduje v podkapitolách.

\begin{table}[h]
\centering
\caption{Prehľad vstupných zdrojov dát aplikácie.}
\label{tab:vstupne-zdroje}
\begin{tabular}{|p{3.2cm}|p{2.6cm}|p{3.2cm}|p{4.2cm}|}
\hline
\textbf{Zdroj} & \textbf{Formát} & \textbf{Prístup} & \textbf{Rola v workflow} \\
\hline
Web of Science & CSV, Excel & manuálny export z webu & primárny vstup metadát \\
\hline
Scopus & CSV & manuálny export z webu & primárny vstup metadát \\
\hline
Repozitár UTB (\texttt{utb\_metadata\_arr}) & PostgreSQL & čítanie a zápis po schválení & cieľové úložisko, referenčný stav pre dedup. \\
\hline
\texttt{utb\_authors\_limited} & PostgreSQL & čítanie & autoritatívny register interných autorov \\
\hline
\texttt{utb\_authors} & PostgreSQL & čítanie a zápis & rozšírený editovateľný register \\
\hline
OBD (\texttt{obd\_publikace}) & PostgreSQL & čítanie & overenie existencie OBDID \\
\hline
Crossref REST API & JSON cez HTTP & čítanie & externé overenie a doplnenie \\
\hline
\end{tabular}
\end{table}

\subsubsection{Externé bibliografické databázy}

Hlavnými vstupnými zdrojmi sú Web of Science a Scopus. Obe platformy umožňujú používateľovi vyexportovať bibliografické záznamy v tabuľkovom formáte, pričom je možné voliť rozsah polí a typ výstupného súboru \cite{ClarivateWOSExport,ScopusSupport}. Knižnica UTB pracuje predovšetkým s exportmi do CSV a Excel a tieto súbory sú v ďalšom workflow chápané ako prvá vrstva surových dát. Záznamy v nich sú formálne štruktúrované, avšak medzi sebou sa líšia v zápise mien autorov, v reprezentácii afiliácií aj v podobe identifikátorov. Pri tom istom diele teda môžu existovať dva mierne odlišné záznamy, jeden z WoS a druhý zo Scopusu, ktoré sa musia v lokálnom prostredí zlúčiť do jedného autoritatívneho záznamu.

\subsubsection{Interný repozitár publikácií UTB}

Druhým zdrojom je samotný interný repozitár knižnice UTB, ktorý je technicky postavený na databáze PostgreSQL. Pre potreby tejto práce sú podstatné najmä tabuľky s metadátovými poľami záznamov, kde sú uložené hodnoty získané z predchádzajúcich importov. Ide o pole, ktoré z jednej strany predstavuje cieľové úložisko (sem sa po schválení knihovníkom zapisujú konečné hodnoty) a z druhej strany aj referenčný stav (existujúce záznamy slúžia pri deduplikácii ako zdroj porovnávania). V repozitári sa používa konvencia stĺpcov v tvare \texttt{dc.*} pre štandardné metadáta podľa Dublin Core a \texttt{utb.*} pre lokálne rozšírenia, napríklad \texttt{utb.wos.affiliation}, \texttt{utb.scopus.affiliation} alebo \texttt{utb.fulltext.dates}.

\subsubsection{Lokálne autoritatívne registre}

Pre identifikáciu interných autorov sa využíva tabuľka \texttt{utb\_authors\_limited}, ktorá obsahuje preddefinované varianty mien zamestnancov UTB v jednotnom kontrolovanom tvare. Okrem nej je k dispozícii rozšírený register \texttt{utb\_authors}, ktorý umožňuje doplňujúcu evidenciu vrátane ďalších identifikátorov. Z dátového hľadiska majú obe tabuľky charakter autoritatívneho záznamu osoby \cite{FRAD2013}: textové pole s menom slúži na vyhľadávanie, no skutočnou referenčnou jednotkou je interný identifikátor osoby. Vstupom pre rozhodnutie, či konkrétny autor záznamu je interným autorom UTB, je zhoda na úrovni týchto registrov, nie len výskyt mena v textovom poli.

Súvisiacim zdrojom je systém OBD, ktorý spravuje organizačnú štruktúru fakúlt, ústavov a pracovísk. V kontexte tejto práce poskytuje OBD aktuálne väzby zamestnanec-pracovisko a slúži ako referenčný rámec pri normalizácii afiliácií na konkrétnu fakultu a ústav.

\subsubsection{Externé otvorené služby}

Doplnkovým zdrojom je Crossref, ktorý poskytuje verejne dostupné REST API s bibliografickými metadátami k DOI, periodikám aj jednotlivým dielam \cite{CrossrefRESTAPI}. V navrhovanom workflow neslúži ako primárny importný kanál, ale ako referenčný zdroj pre overenie a doplnenie hodnôt, najmä pri normalizácii časopisov a pri kontrole formálnej správnosti DOI. Crossref zároveň deklaruje, že väčšina jeho bibliografických metadát je voľne použiteľná \cite{CrossrefMetadataRetrieval}, čo znižuje právne riziká pri jeho zaradení do lokálneho spracovania.

\subsection{Analýza pôvodného workflow knižnice UTB}

Pred zavedením navrhovanej aplikácie pracovala knižnica UTB s publikačnými záznamami v postupnosti, ktorú možno rozdeliť do troch krokov. V prvom kroku boli záznamy vyexportované zo Scopusu a Web of Science do tabuľkového formátu. V druhom kroku prebiehala technická kontrola hrubých chýb a importérske spracovanie. V treťom kroku pracovníci knižnice v prostredí Google Sheets ručne dopĺňali a opravovali polia, ktoré boli neúplné, nekonzistentné alebo nezodpovedali interným pravidlám evidencie. Až po tomto kroku sa údaje ukladali do internej databázy. Tento postup je dlhodobo pragmatický, no jeho slabiny sú zrejmé z empirických pozorovaní pri spolupráci s pracovníkmi knižnice a zo systematickej analýzy príbuzných workflow \cite{Mikulecky2025}.

Z opakovaných pozorovaní vychádzajú štyri kategórie problémov. Prvou je menová nejednoznačnosť autorov: tá istá osoba sa môže v rôznych zdrojoch objaviť ako „Belas, Jaroslav", „Belás, Jaroslav" alebo so skrátenou iniciálou. Druhou je formátovanie dátumov, kde sa popri ISO formátoch vyskytujú aj americké zápisy v tvare MDR alebo voľný text v poli \texttt{utb.fulltext.dates}. Treťou je nejednotné označovanie časopisov, kde sa rovnaký časopis objavuje s plným aj skráteným názvom, prípadne bez zhody v ISSN. Štvrtou kategóriou sú duplicity vznikajúce kombináciou exportov z WoS a zo Scopusu, kde dva záznamy popisujú tú istú publikáciu, no líšia sa v drobnostiach.

Spoločnou črtou týchto problémov je, že ich nemožno vyriešiť jediným pravidlom. Časť opráv je formálna a deterministická, časť vyžaduje porovnanie s externým zdrojom a časť je interpretačná, pretože medzi dvomi rôzne zapísanými hodnotami treba rozhodnúť, či označujú tú istú entitu, alebo nie. Dôsledkom je, že manuálne spracovanie v Sheets je časovo náročné, repetitívne a slabo škálovateľné, pričom súčasne nemá podporu pre auditovateľnosť: pri pohľade na finálny záznam nie je možné spätne určiť, či konkrétna hodnota pochádza priamo z exportu, alebo bola doplnená rozhodnutím knihovníka.

Nejednotnosť reprezentácie ilustruje konkrétny záznam jednej publikácie tak, ako sa objavil v exporte z Web of Science, v exporte zo Scopusu a po normalizácii v repozitári UTB (Obr.~\ref{fig:author-variants}). Z príkladu je vidno, že každý zo zdrojov ukladá rovnaké informácie iným spôsobom. Web of Science zoskupuje autorov priamo s ich afiliáciou do jedného poľa pomocou hranatých zátvoriek, Scopus eviduje len afiliácie ako samostatný zoznam bez priameho previazania na konkrétne meno autora a repozitár drží mená autorov v normalizovanej podobe oddelené reťazcom \texttt{||}. K štruktúrnym rozdielom sa pripája aj odlišné nakladanie s diakritikou: oba exporty ju v menách aj v názvoch inštitúcií strácajú, kým repozitár ju zachováva vrátane špeciálnych znakov, napríklad mena \texttt{Çera}. Tento jeden záznam tak zachytáva menovú nejednoznačnosť autorov a zároveň ukazuje, prečo aplikácia musí vedieť zladiť tie isté údaje uložené v rôznych poliach a tvaroch naprieč zdrojmi.

\begin{figure}[h]
\centering
\begin{minipage}{0.95\textwidth}

\textbf{Web of Science (\texttt{utb.wos.affiliation})}
\begin{verbatim}
[Belas, Jaroslav; Dvorsky, Jan; Cera, Gentjan]
   Tomas Bata Univ Zlin, Mostni 5139, Zlin 76001, Czech Republic;
[Strnad, Zdenek]
   Univ South Bohemia Ceske Budejovice, Branisovska 31a,
   Ceske Budejovice 37005, Czech Republic;
[Valaskova, Katarina]
   Univ Zilina, Univ 8215-1, Zilina 01026, Slovakia
\end{verbatim}

\textbf{Scopus (pole \texttt{Authors with affiliations})}
\begin{verbatim}
Tomas Bata University in Zlin, Mostni 5139,
   Zlin, 760 01, Czech Republic;
University of South Bohemia in Ceske Budejovice,
   Branisovska 31a, Ceske Budejovice, 370 05, Czech Republic;
University of Zilina, Univerzitna 8215/1,
   Zilina, 010 26, Slovakia
\end{verbatim}

\textbf{Repozitár UTB po normalizácii (\texttt{dc.contributor.author})}
\begin{verbatim}
Belás, Jaroslav||Dvorský, Ján||Strnad, Zdeněk||
   Valášková, Katarína||Çera, Gentjan
\end{verbatim}

\end{minipage}
\caption{Reprezentácia tej istej publikácie v troch zdrojoch. WoS zoskupuje autorov a ich afiliácie do jedného poľa s hranatými zátvorkami, Scopus eviduje len zoznam afiliácií bez previazania na konkrétne meno autora a repozitár drží normalizované mená autorov. K štruktúrnym rozdielom sa pripája chýbajúca diakritika v oboch exportoch.}
\label{fig:author-variants}
\end{figure}

\subsection{Funkčné požiadavky}

Z predchádzajúcej analýzy vyplývajú nasledujúce funkčné požiadavky. Sú očíslované, aby bolo možné v ďalších kapitolách jednoznačne odkazovať na ich realizáciu a vyhodnotenie.

\begin{itemize}
    \item \textbf{FR1: Import záznamov z WoS a Scopus.} Aplikácia musí načítať bibliografické záznamy z exportov WoS a Scopus do lokálnej pracovnej databázy a uchovať ich vo forme, ktorá zachová pôvodné hodnoty pre neskoršie porovnanie.
    \item \textbf{FR2: Validácia a data hygiene.} Aplikácia musí identifikovať bežné formálne chyby (nepovolené znaky, mojibake, nadbytočné medzery, nesprávne sformátované identifikátory) a navrhnúť ich opravy oddelene od skutočného zápisu hodnôt.
    \item \textbf{FR3: Detekcia interných autorov UTB.} Aplikácia musí porovnať autorov záznamu voči autoritatívnemu registru UTB, identifikovať podmnožinu interných autorov a priradiť im fakultu a ústav podľa afiliácie.
    \item \textbf{FR4: Normalizácia dátumov.} Aplikácia musí extrahovať redakčné dátumy z poľa \texttt{utb.fulltext.dates} a z ďalších zdrojov a previesť ich do formátu ISO 8601 \cite{ISO8601_2019}, vrátane logickej kontroly chronologického poradia.
    \item \textbf{FR5: Normalizácia časopisov a vydavateľov.} Aplikácia musí mapovať názov časopisu a vydavateľa na konzistentnú podobu, prednostne podľa ISSN, ISBN alebo DOI, s možnosťou doplnenia z Crossref \cite{CrossrefRESTAPI,ISSNWhatIs}.
    \item \textbf{FR6: Deduplikácia záznamov.} Aplikácia musí odhaliť duplicitné záznamy v rámci jedného importu, medzi importnými dávkami a voči existujúcim záznamom repozitára, primárne podľa DOI a sekundárne podľa kombinácie atribútov.
    \item \textbf{FR7: Asistencia LLM tam, kde heuristiky nestačia.} V krokoch detekcie autorov a extrakcie dátumov musí byť k dispozícii voliteľný LLM fallback so štruktúrovaným výstupom v JSON, ktorý sa použije len pri prípadoch, ktoré deterministické pravidlá nedokázali jednoznačne vyriešiť.
    \item \textbf{FR8: Asistovaná editácia záznamu knihovníkom.} Aplikácia musí poskytnúť používateľské rozhranie, v ktorom knihovník vidí porovnanie hodnôt z viacerých zdrojov, môže ich upravovať a manuálne potvrdiť výslednú podobu záznamu.
    \item \textbf{FR9: Schvaľovací zásobník zmien.} Úpravy a navrhnuté hodnoty sa pred zápisom do produkčnej databázy uchovávajú v zásobníku, kde ich knihovník môže preverovať, vracať a finálne schvaľovať.
    \item \textbf{FR10: Audit pôvodu hodnôt.} Pri každej hodnote musí byť možné určiť, či pochádza z pôvodného importu, z heuristického návrhu, z výstupu LLM alebo z manuálneho zásahu knihovníka.
\end{itemize}

\subsection{Nefunkčné požiadavky}

Vedľa funkčných požiadaviek bolo počas analýzy a konzultácií so zadávateľom sformulovaných niekoľko nefunkčných obmedzení a očakávaní, ktoré ovplyvňujú návrh systému.

\begin{itemize}
    \item \textbf{NFR1: Local-first prevádzka.} Pri rutinných úlohách musí byť možné spustiť celé spracovanie lokálne, bez závislosti od externého LLM API, aby sa minimalizoval tok citlivých údajov mimo organizácie a aby bola znížená závislosť od poskytovateľa.
    \item \textbf{NFR2: Vymeniteľnosť LLM vrstvy.} Aplikačný kód nesmie byť pevne viazaný na konkrétneho LLM poskytovateľa. Cez jednotné rozhranie musí byť možné prepnúť medzi lokálnym (Ollama) a cloudovým (OpenAI-kompatibilné API) režimom.
    \item \textbf{NFR3: Opakovateľnosť spracovania.} Pipeline musí byť možné spustiť opakovane, idempotentne, bez toho, aby opakovaný beh skreslil už raz správne odvodené hodnoty.
    \item \textbf{NFR4: Striktný JSON výstup z LLM.} Odpovede modelu sa musia podriaďovať preddefinovanej schéme tak, aby ich bolo možné automaticky validovať \cite{OpenAIStructuredOutputs,OllamaStructuredOutputs}.
    \item \textbf{NFR5: Auditovateľnosť úprav.} Každý zásah do záznamu musí byť rozlíšiteľný podľa pôvodu (heuristika, LLM, človek), čo je v súlade s konceptom provenancie údajov \cite{PROVOverview}.
    \item \textbf{NFR6: Kontrolovaný prístup.} Aplikácia má byť dostupná iba pre vymenovaných pracovníkov knižnice, s prihlásením integrovaným do prostredia UTB \cite{Konarik2025}.
    \item \textbf{NFR7: Externé plánovanie behov.} Periodické spustenie pipeline má byť riešené mimo aplikácie, prostredníctvom Windows Task Scheduler alebo cron, aby sa aplikácia nemusela starať o dlhodobo bežiace plánovacie procesy.
\end{itemize}

\subsection{Scope a vymedzenie hraníc}

Pre korektnú obhajobu navrhovaného riešenia je dôležité jednoznačne pomenovať aj to, čo aplikácia robiť \emph{nebude}. Bez takého vymedzenia by sa od nej dali očakávať vlastnosti, ktoré nikdy neboli súčasťou zadania.

V rámci tejto práce je predmetom realizácie nasledujúce: spracovanie a normalizácia bibliografických metadát z vopred získaných exportov WoS a Scopus, ich konsolidácia s existujúcim repozitárom knižnice, asistovaná editácia záznamov a uloženie schválených hodnôt v rámci kontrolovaného workflow. Aplikácia podporuje pracovníka knižnice v rutinných opravách a normalizáciách, no neodstraňuje nutnosť ľudského schvaľovania pri rozhodovaní o identite osôb, časopisov a duplicít.

Predmetom práce naopak \emph{nie je} nasledujúce: nahradenie existujúceho repozitára publikácií UTB, automatizovaný zber dát z webu mimo dohodnutých zdrojov, plnotextová indexácia článkov, vlastná autentifikačná infraštruktúra mimo už zavedeného UTB prihlasovania ani trénovanie alebo doladenie samotných jazykových modelov. Aplikácia tiež nepretláča lokálne autoritatívne záznamy do externých systémov a nie je nástrojom pre koncového používateľa repozitára (čitateľa, autora). Jej cieľovou skupinou sú výlučne pracovníci knižnice, ktorí spravujú evidenciu publikačnej činnosti.

Tým je rámec praktickej časti vymedzený tak, aby zodpovedal zadaniu \emph{Aplikácia pre podporu správy repozitáre vedeckých výstupov pomocou AI}, no zároveň zostal realisticky obhájiteľný v termíne odovzdania a počas obhajoby.
% Kapitola 8: Návrh technického riešenia
% Stav: draft

\section{Návrh technického riešenia}

Predchádzajúca kapitola pomenovala vstupné zdroje, problémy súčasného workflow a požiadavky, ktoré z neho vyplývajú. Na týchto východiskách stojí návrh technického riešenia. Cieľom tejto kapitoly je opísať architektúru aplikácie, jej dátový model a kľúčové návrhové rozhodnutia. Dôraz je kladený na to, aby boli jednotlivé časti vzájomne oddeliteľné, aby výstupy automatických krokov bolo možné spätne preskúmať a aby finálne rozhodnutie nad záznamom zostalo v rukách knihovníka.

\subsection{Architektúra na úrovni komponentov}

Aplikácia je rozdelená do dvoch logických vrstiev. Prvou je dátový pipeline, ktorý beží ako sada príkazov v príkazovom riadku a vykonáva nemiešanú prácu nad dátami, teda import, validáciu, detekciu autorov, normalizáciu dátumov a časopisov a deduplikáciu. Druhou vrstvou je webová aplikácia napísaná vo Flasku, ktorá poskytuje knihovníkovi rozhranie pre kontrolu, úpravu a schvaľovanie záznamov. Obe vrstvy zdieľajú spoločnú lokálnu databázu PostgreSQL a knižnicu spoločných modulov v adresári \texttt{src/}.

Toto rozdelenie sleduje dva ciele. Po prvé umožňuje spúšťať dávkové spracovanie nezávisle od webovej aplikácie, napríklad v noci alebo v naplánovaných intervaloch cez systémový plánovač. Po druhé oddeľuje výpočtovo náročné kroky od krokov, ktoré vyžadujú interakciu s človekom. Pipeline pripravuje návrhy, používateľské rozhranie tieto návrhy vizualizuje a knihovník ich schvaľuje alebo upravuje.

Z hľadiska tokov dát aplikácia komunikuje s tromi okruhmi systémov. Vstupom je interný repozitár publikácií knižnice UTB, ku ktorému sa pristupuje len na čítanie, prípadne zápis až vo finálnom kroku schválenia. Druhým okruhom sú externé bibliografické databázy WoS a Scopus, ktorých exporty sa importujú lokálne. Tretím okruhom sú externé referenčné služby, najmä Crossref REST API \cite{CrossrefRESTAPI}, voči ktorým aplikácia overuje a dopĺňa metadáta. Všetky LLM volania prebiehajú cez jednotný adaptér, ktorý sa dá v konfigurácii prepnúť medzi lokálnym a cloudovým režimom.

Pre potreby tejto kapitoly stačí poznamenať, že každý krok pipeline zapisuje výstup do oddelenej pracovnej tabuľky, nie priamo do produkčného repozitára. Schematické zachytenie celkového dátového toku ako samostatný diagram je vhodné doplniť pri finálnej úprave práce.
% TODO: doplniť diagram architektúry na úrovni komponentov; vložiť ako figure s label fig:architektura a v texte naň odkázať.

\subsection{Dátový model a pracovné tabuľky}

Centrálnou myšlienkou dátového modelu je oddelenie produkčných údajov, pracovných výstupov pipeline a zásobníka schvaľovaných zmien do troch samostatných vrstiev. Vďaka tomu je možné automatické kroky opakovane spúšťať bez rizika, že prepíšu manuálne potvrdené hodnoty.

Prvú vrstvu tvorí lokálna kópia produkčnej tabuľky repozitára (\texttt{utb\_metadata\_arr}), ktorá vzniká príkazom \texttt{bootstrap-local-db}. Tento krok skopíruje stĺpce z remote PostgreSQL inštancie do lokálnej databázy a slúži ako vstupný snapshot pre celé spracovanie. Implementácia kopírovania je v \href{src/db/setup.py}{src/db/setup.py}, kde sa metadáta stĺpcov najprv načítajú z \texttt{information\_schema.columns} a následne sa použijú na vygenerovanie zodpovedajúceho \texttt{CREATE TABLE} v lokálnej databáze.

Druhú vrstvu predstavuje pracovná tabuľka \texttt{utb\_processing\_queue}, do ktorej zapisujú všetky pipeline kroky. Štruktúru tabuľky a sadu výstupných stĺpcov definuje konštanta \texttt{OUTPUT\_COLUMNS} v module \href{src/common/constants.py}{src/common/constants.py}. Každá fáza pipeline má v tejto tabuľke vyhradené stĺpce pre svoj výstup (napríklad \texttt{author\_internal\_names}, \texttt{author\_faculty}, \texttt{utb\_date\_received}), pre stav spracovania (\texttt{author\_heuristic\_status}, \texttt{author\_llm\_status}) a pre príznaky (\texttt{author\_flags}, \texttt{date\_flags}), kam sa uchovávajú metadáta o priebehu rozhodnutia. Stavy sú reprezentované cez triedy \texttt{HeuristicStatus}, \texttt{LLMStatus}, \texttt{ValidationStatus} a \texttt{DateLLMStatus}, čím je zabezpečená konzistencia hodnôt v celej aplikácii.

Treťou vrstvou je zásobník zmien \texttt{utb\_change\_buffer}, kde sa uchovávajú úpravy navrhnuté knihovníkom. Toto je miesto, kde končí webová aplikácia, keď používateľ stlačí Save alebo \texttt{Ctrl+S}. Až po explicitnom schválení sa zmeny premietajú do produkčnej databázy. Súčasťou modelu je aj samostatná história deduplikácie \texttt{dedup\_history}, ktorá zaznamenáva, ktoré záznamy boli zlúčené s ktorými a na základe akých kľúčov.

Pre lokálne autoritatívne registre aplikácia využíva tabuľky \texttt{utb\_authors\_limited} a \texttt{utb\_authors}. Prvá obsahuje varianty mien interných autorov UTB a slúži najmä pri detekcii. Druhá je rozšírený editovateľný register, do ktorého môže knihovník pridávať nové osoby. Pre validáciu OBDID identifikátorov sa pristupuje na čítanie do remote tabuľky \texttt{veda.obd\_publikace}, čo je popísané v module \href{src/quality/checks.py}{src/quality/checks.py}, konkrétne vo funkcii \texttt{check\_obdid\_batch}.

\subsection{Pipeline fázy a ich postupnosť}

Pipeline je navrhnutý tak, aby každú fázu bolo možné spustiť samostatne. Súčasne má odporúčané poradie, v ktorom výstup jednej fázy vstupuje ako kontext do ďalšej. Inicializačné kroky pripravia tabuľky a skopírujú dáta. Validácia označí formálne chyby a navrhne ich opravy. Detekcia autorov pracuje s tým, čo prišlo z importu alebo z validácie. Extrakcia dátumov a normalizácia časopisov sú navzájom nezávislé a možno ich vykonávať v ľubovoľnom poradí. Deduplikácia je na konci, keď sú metadáta v očakávanej podobe.

Rozdelenie do CLI príkazov, ktoré toto poradie dodržuje, je nasledovné:

\begin{itemize}
    \item \texttt{bootstrap-local-db}, \texttt{setup-processing-queue}, \texttt{setup-dedup-history}: inicializácia tabuliek a kópia zo zdroja.
    \item \texttt{validate-metadata}, \texttt{apply-validation-fixes}, \texttt{metadata-validation-status}: kontrola formálnych chýb a aplikácia návrhov.
    \item \texttt{detect-authors}, \texttt{detect-authors-llm}, \texttt{compare-author-detection}, \texttt{author-detection-status}: heuristická detekcia interných autorov a voliteľný LLM fallback.
    \item \texttt{extract-dates}, \texttt{extract-dates-llm}, \texttt{date-extraction-status}: parsovanie redakčných dátumov.
    \item \texttt{normalize-journals}, \texttt{apply-journal-normalization}, \texttt{journal-normalization-status}: normalizácia časopisov a vydavateľov.
    \item \texttt{deduplicate-records}, \texttt{deduplication-status}: odhalenie a zlúčenie duplicít.
\end{itemize}

LLM fallbacky sú zámerne samostatné príkazy, nie automatické pokračovanie heuristík. Dôvodom je, že lokálny aj cloudový model majú vyššiu prevádzkovú cenu (čas, prípadne API náklady) a v praxi sa spúšťajú len pre záznamy, ktoré heuristika označila ako nejednoznačné, teda nastavila stĺpec \texttt{author\_needs\_llm} alebo \texttt{date\_needs\_llm} na \texttt{TRUE}.

Celý katalóg príkazov a ich popis pre webové rozhranie definuje \href{web/blueprints/pipeline/catalog.py}{web/blueprints/pipeline/catalog.py}. Skutočné spustenie príkazu z UI obsluhuje SSE runner v \href{web/blueprints/pipeline/routes.py}{web/blueprints/pipeline/routes.py}, ktorý streamuje výstup procesu do prehliadača v reálnom čase.

\subsection{LLM komponenta}

Použitie jazykového modelu je v aplikácii navrhnuté ako vymeniteľná vrstva s jednotným rozhraním. Cieľom je, aby aplikačný kód, ktorý LLM volá, nemusel vedieť, či bežne odpovedá lokálne nasadený Ollama model, alebo cloudový model za OpenAI-kompatibilným API. Toto rozhranie je v \href{src/llm/client.py}{src/llm/client.py}, kde sú implementované triedy \texttt{OllamaClient} a \texttt{CloudLLMCompatibleClient}, obe odvodené od abstraktnej \texttt{LLMClient}. Výber poskytovateľa určuje konfiguračná premenná \texttt{LLM\_PROVIDER}, ktorá sa načítava v \href{src/config/settings.py}{src/config/settings.py}.

Nad klientom je postavená trieda \texttt{LLMSession} z modulu \href{src/llm/session.py}{src/llm/session.py}, ktorá zapuzdruje fixný kontext spracovania (systémový prompt, JSON schému a voliteľnú preambulu) a poskytuje jednoduché \texttt{ask(user\_message)} volanie pre každý spracúvaný záznam. Pre Ollamu sa využíva preamble vo forme úvodného dialógu s ukážkou správneho výstupu. Tento dialóg je rovnaký pre celú dávku záznamov, čo umožňuje modelu zdieľať predpočítané KV-cache medzi volaniami a urýchľuje spracovanie. Pre cloudový režim sa preamble nepoužíva, lebo zdieľanie KV-cache nie je v štandardnom OpenAI API exponované rovnakým spôsobom.

Štruktúrované výstupy sú vynútené cez JSON schému, ktorá je súčasťou každej úlohy. Pre Ollamu sa schéma prenáša cez parameter \texttt{format} \cite{OllamaStructuredOutputs}, pre cloudový režim aplikácia využíva \texttt{response\_format} alebo function calling podľa toho, čo poskytovateľ podporuje \cite{OpenAIStructuredOutputs,OpenAIFunctionCalling}. Konkrétne úlohy (extrakcia interných autorov, parsovanie dátumov) sú implementované ako samostatné moduly v \href{src/llm/tasks/}{src/llm/tasks/}, kde každý modul definuje svoj systémový prompt, schému a parser výstupu.

\subsection{Zdôvodnenie výberu technológií}

Pri výbere konkrétnych nástrojov bola hlavným kritériom dlhodobá udržateľnosť riešenia v prostredí univerzitnej knižnice, kompatibilita s existujúcou infraštruktúrou a dostupnosť skúseností v rámci predchádzajúcich diplomových prác na FAI UTB.

Ako programovací jazyk bol zvolený Python verzie 3.11 a vyšší. Dôvodom je široká dostupnosť knižníc pre prácu s dátami, integrácie s LLM klientmi a fakt, že rovnaký jazyk používa aj príbuzná diplomová práca P. Mikuleckého \cite{Mikulecky2025}, čo uľahčuje opätovné využitie myšlienok a porovnanie prístupov. Pre CLI rozhranie bol zvolený Typer, ktorý poskytuje deklaratívny zápis príkazov a automatické generovanie nápovedy. Webová vrstva je postavená na Flasku, čo nadväzuje na bakalársku prácu M. Koňaříka \cite{Konarik2025} a umožňuje prebrať existujúce vzory pre autentifikáciu UTB a štruktúru blueprintov. Plnohodnotnejší framework typu Django by pri rozsahu tejto aplikácie znamenal viac réžie, ako prínosov.

Databázová vrstva je realizovaná cez SQLAlchemy 2.x s drajverom psycopg verzie 3 a pripája sa k PostgreSQL. Voľba PostgreSQL je daná tým, že na ňom beží aj produkčný repozitár UTB a aplikácia tak môže pracovať s rovnakými dátovými typmi vrátane polí (\texttt{TEXT[]}) a JSONB. Pre validáciu Python objektov sa používa Pydantic verzie 2 a pre parsovanie nečistých textov knižnica \texttt{ftfy}, ktorá rieši typické prípady mojibake. Pre fuzzy porovnávanie reťazcov je nasadená knižnica \texttt{jellyfish}, ktorá pokrýva edit-distance metriky aj Jaro-Winklerovu podobnosť \cite{Christen2012,Winkler1990}. HTTP volania na Crossref API obsluhuje \texttt{httpx} s podporou opakovaných pokusov a časových limitov.

LLM vrstva má dve referenčné nasadenia. Pre lokálne behy sa používa Ollama s modelom radu \texttt{qwen3} v menšej veľkosti (8B parametrov), ktorá sa zmestí na bežné kancelárske GPU alebo aj na CPU pri rozumných časoch odozvy. Pre cloudové behy je ako referenčný poskytovateľ zvolený Groq, ktorého API explicitne dokumentuje OpenAI-kompatibilitu, takže ho možno volať tým istým klientom ako oficiálne OpenAI \cite{GroqOpenAICompatibility}. Aplikácia teda nezávisí od konkrétneho cloudového poskytovateľa a vie sa preorientovať podľa potreby.

Pri prieskumnej analýze dát počas vývoja sa využila aj utilita DuckDB, najmä pri overovaní hypotéz nad exportmi WoS a Scopus a pri rýchlej kontrole, čo presne sa nachádza v poli \texttt{utb.fulltext.dates} v jednotlivých rokoch. DuckDB nie je súčasťou bežne behajúcej aplikácie, no v práci je relevantný preto, lebo formoval rozhodnutia o tom, ktoré dátové pravidlá sa dali pokryť deterministicky a kde bol naozaj potrebný LLM fallback.

\subsection{Wireframy a UI návrh}

Návrh používateľského rozhrania vychádzal z toho, ako knihovník reálne pracuje s tabuľkou v Google Sheets. Rozhranie preto kopíruje stĺpcový pohľad, kde popri sebe stoja hodnoty z rôznych zdrojov (Repozitár, WoS, Scopus, Crossref) a knihovník môže vizuálne porovnávať rozdiely. Cieľom nie je nahradiť ho úplne odlišnou interakciou, ale doplniť mu nástroje, ktoré v Sheets chýbali, najmä zvýrazňovanie rozdielov, asistované návrhy a auditovateľný zásobník zmien.

Aplikácia má tri hlavné obrazovky:

\begin{itemize}
    \item \textbf{Zoznam záznamov} (\texttt{/}), kde knihovník vidí prehľad importovaných publikácií, ich stav spracovania, prípadné varovania a duplicitné dvojice.
    \item \textbf{Detail záznamu} (\texttt{/record/<resource\_id>}), ktorý je centrálnou pracovnou plochou. Obsahuje editovateľnú tabuľku polí so zvýraznenými rozdielmi medzi zdrojmi, akcie pre uloženie zmien do zásobníka a krátku históriu.
    \item \textbf{Nastavenia pipeline} (\texttt{/settings/pipeline}), kde je možné spustiť jednotlivé príkazy pipeline aj v rámci odporúčaného poradia. Stránka prijíma SSE výstup z bežiaceho procesu a zobrazuje ho v reálnom čase.
\end{itemize}

Pomocné nastavenia, napríklad konfigurovateľné poradie riadkov v detaile záznamu (\texttt{/settings/row-order}), sú k dispozícii pre pracovníkov, ktorí preferujú odlišné usporiadanie polí. Wireframy detailu záznamu, vrátane editácie autorov a afiliácií, vznikli ešte pred implementáciou v podobe interaktívnych prototypov a vychádzali z reálneho príkladu duplicitného záznamu z Web of Science a Scopusu.



% ============================================================================ %
% Encoding: UTF-8 (žluťoučký kůň úpěl ďábelšké ódy)
% ============================================================================ %

\sectionwithoutnumbering{Závěr}

Text \dots


% ============================================================================ %


\paragraphindentationstop

% ============================================================================ %

% Využívají se bibliografické údaje uložené v souboru tex/references.bib.
% Ty se automaticky upravuji dle zvolené citační normy (ČSN ISO 690:2022).
\insertlistofreferences

%% ============================================================================ %
% Encoding: UTF-8 (žluťoučký kůň úpěl ďábelšké ódy)
% ============================================================================ %

% ---------------------------------------------------------------------------- %
\addappendix{První příloha}
Obsah přílohy


% ---------------------------------------------------------------------------- %
\addappendix{Druhá příloha}
Obsah přílohy

\addappendixother{Označení}{Název přílohy}{Soubor}

% ============================================================================ %
  % Vložení příloh

% ============================================================================ %

\end{document}

% ============================================================================ %
